\documentclass[10pt,twocolumn]{article}

% ------------------------------------------------------------
% Page layout: compact two-column proceedings-style layout
% ------------------------------------------------------------
\usepackage[
  letterpaper,
  top=0.75in,
  bottom=1.0in,
  left=0.625in,
  right=0.625in
]{geometry}

\usepackage{placeins}

\setlength{\columnsep}{0.25in}

\usepackage{microtype}
\usepackage{xurl}
\usepackage{seqsplit}
\setlength{\emergencystretch}{2em}
\hbadness=10000
\vbadness=10000
\hfuzz=20pt
\vfuzz=20pt
\frenchspacing

% Compact paragraph style: no blank paragraph gaps.
\setlength{\parindent}{1em}
\setlength{\parskip}{0pt}

% Slightly tighter line spacing.
\linespread{0.97}

% ------------------------------------------------------------
% Core math
% Load amsmath/mathtools before unicode-math.
% ------------------------------------------------------------
\usepackage{amsmath}
\usepackage{mathtools}

% ------------------------------------------------------------
% Fonts: LuaLaTeX or XeLaTeX
% ------------------------------------------------------------
\usepackage{fontspec}
\usepackage[warnings-off={mathtools-colon,mathtools-overbracket}]{unicode-math}

% Use Times-like fonts if available; otherwise fall back to DejaVu.
\IfFontExistsTF{Times New Roman}
  {\setmainfont{Times New Roman}}
  {%
    \IfFontExistsTF{TeX Gyre Termes}
      {\setmainfont{TeX Gyre Termes}}
      {\setmainfont{DejaVu Serif}}%
  }
  
\IfFontExistsTF{TeX Gyre Termes}
  {\setmainfont{TeX Gyre Termes}}
  {\setmainfont{DejaVu Serif}}

\IfFontExistsTF{TeX Gyre Heros}
  {\setsansfont{TeX Gyre Heros}}
  {\setsansfont{DejaVu Sans}}

\setmonofont{DejaVu Sans Mono}[Scale=0.78]

\IfFontExistsTF{TeX Gyre Termes Math}
  {\setmathfont{TeX Gyre Termes Math}}
  {\setmathfont{Latin Modern Math}}

% Do NOT load amssymb, amsfonts, or bm with unicode-math.
% amssymb/amfonts can conflict with symbols such as \eth.
% bm can break Greek bold symbols under unicode-math.

% ------------------------------------------------------------
% Extra math utilities
% ------------------------------------------------------------
\usepackage{tensor}
\usepackage{cancel}

% ------------------------------------------------------------
% Formal compact section headings
% ------------------------------------------------------------
\usepackage{titlesec}

% Section numbering:
% I. Section
% A. Subsection
% 1) Subsubsection
\renewcommand{\thesection}{\Roman{section}}
\renewcommand{\thesubsection}{\Alph{subsection}}
\renewcommand{\thesubsubsection}{\arabic{subsubsection}}

\titleformat{\section}
  [block]
  {\normalfont\normalsize\scshape\centering}
  {\thesection.}
  {0.5em}
  {}

\titleformat{\subsection}
  [block]
  {\normalfont\normalsize\itshape}
  {\thesubsection.}
  {0.5em}
  {}

\titleformat{\subsubsection}
  [runin]
  {\normalfont\normalsize\itshape}
  {\thesubsubsection)}
  {0.5em}
  {}
  [. ]

\titlespacing*{\section}
  {0pt}
  {1.4ex plus 0.4ex minus 0.2ex}
  {0.7ex}

\titlespacing*{\subsection}
  {0pt}
  {1.0ex plus 0.3ex minus 0.2ex}
  {0.5ex}

\titlespacing*{\subsubsection}
  {\parindent}
  {0.8ex plus 0.2ex minus 0.2ex}
  {0.5em}

% ------------------------------------------------------------
% Compact lists
% ------------------------------------------------------------
\usepackage{enumitem}

\setlist{
  nosep,
  leftmargin=1.2em
}

% ------------------------------------------------------------
% Tables
% ------------------------------------------------------------
\usepackage{booktabs}
\usepackage{array}
\usepackage{tabularx}
\usepackage{multirow}
\usepackage{siunitx}
\usepackage{ragged2e}
\usepackage{etoolbox}

\sisetup{
  detect-all,
  per-mode=symbol,
  exponent-product=\cdot
}

\setlength{\tabcolsep}{3.2pt}
\renewcommand{\arraystretch}{1.03}

\newcolumntype{Y}{>{\RaggedRight\arraybackslash\hspace{0pt}}X}
\newcolumntype{Z}[1]{>{\RaggedRight\arraybackslash\hspace{0pt}}p{#1}}

% Default table text size. Use {\scriptsize ...} locally for very wide tables.
\AtBeginEnvironment{tabular}{\footnotesize}
\AtBeginEnvironment{tabularx}{\footnotesize}

% ------------------------------------------------------------
% Figures, SVG, PGF/TikZ plots
% ------------------------------------------------------------
\usepackage{graphicx}
\usepackage{adjustbox}
\usepackage{svg}
\usepackage{tikz}
\usepackage{pgfplots}
\usepackage{stfloats}

\pgfplotsset{compat=1.18}

\usetikzlibrary{
  arrows.meta,
  calc,
  positioning,
  decorations.pathmorphing,
  matrix
}

% ------------------------------------------------------------
% Compact floats and equations
% ------------------------------------------------------------

% Allow LaTeX to place more floats before deferring them to the end.
\renewcommand{\topfraction}{0.95}
\renewcommand{\bottomfraction}{0.90}
\renewcommand{\textfraction}{0.05}
\renewcommand{\floatpagefraction}{0.80}

\renewcommand{\dbltopfraction}{0.95}
\renewcommand{\dblfloatpagefraction}{0.80}

\setcounter{topnumber}{5}
\setcounter{bottomnumber}{5}
\setcounter{totalnumber}{10}
\setcounter{dbltopnumber}{5}

% Normal float spacing.
\setlength{\floatsep}{7pt plus 1pt minus 1pt}
\setlength{\textfloatsep}{8pt plus 1pt minus 1pt}
\setlength{\intextsep}{7pt plus 1pt minus 1pt}

% Double-column float spacing.
\setlength{\dblfloatsep}{7pt plus 1pt minus 1pt}
\setlength{\dbltextfloatsep}{8pt plus 1pt minus 1pt}

% Prevent float-only pages from spreading tables vertically.
\makeatletter
\setlength{\@fptop}{0pt}
\setlength{\@fpsep}{7pt}
\setlength{\@fpbot}{0pt plus 1fil}

\setlength{\@dblfptop}{0pt}
\setlength{\@dblfpsep}{7pt}
\setlength{\@dblfpbot}{0pt plus 1fil}
\makeatother

\AtBeginDocument{%
  \setlength{\abovedisplayskip}{6pt plus 2pt minus 2pt}%
  \setlength{\belowdisplayskip}{6pt plus 2pt minus 2pt}%
  \setlength{\abovedisplayshortskip}{4pt plus 2pt minus 2pt}%
  \setlength{\belowdisplayshortskip}{4pt plus 2pt minus 2pt}%
}

% ------------------------------------------------------------
% Captions and references
% ------------------------------------------------------------
\usepackage{caption}
\usepackage{subcaption}

\captionsetup{
  font=footnotesize,
  labelfont=bf,
  skip=3pt
}

\captionsetup[table]{
  position=top,
  font=footnotesize,
  labelfont=bf,
  skip=3pt
}

\captionsetup[figure]{
  font=footnotesize,
  labelfont=bf,
  skip=3pt
}

\usepackage[hidelinks]{hyperref}
\usepackage[nameinlink,noabbrev]{cleveref}

% ------------------------------------------------------------
% Compact title block
% ------------------------------------------------------------
\makeatletter
\renewcommand{\maketitle}{%
  \twocolumn[{%
    \vspace*{-0.35in}
    \begin{center}
      {\normalfont\rmfamily\fontsize{24}{28}\selectfont\@title\par}
      \vspace{0.9em}
      {\normalfont\normalsize\rmfamily\@author\par}
      \ifx\@date\@empty
        \vspace{0.7em}
      \else
        \vspace{0.4em}
        {\normalfont\normalsize\rmfamily\@date\par}
        \vspace{0.7em}
      \fi
    \end{center}
  }]%
}
\makeatother

% ------------------------------------------------------------
% Abstract and keywords style
% ------------------------------------------------------------
\renewenvironment{abstract}
  {%
    \par
    \small
    \bfseries
    \noindent
    \textit{Abstract---}%
    \ignorespaces
  }
  {%
    \par
    \vspace{0.7em}
  }

\newenvironment{keywords}
  {%
    \par
    \small
    \noindent
    \textbf{\textit{Keywords---}}%
    \ignorespaces
  }
  {%
    \par
    \vspace{1.0em}
  }

% ------------------------------------------------------------
% Common math macros
% ------------------------------------------------------------

% Sets and spaces
\newcommand{\R}{\mathbb{R}}
\newcommand{\C}{\mathbb{C}}
\newcommand{\N}{\mathbb{N}}
\newcommand{\Z}{\mathbb{Z}}
\newcommand{\Q}{\mathbb{Q}}

\newcommand{\Ltwo}{L^2}
\newcommand{\Lp}[1]{L^{#1}}
\newcommand{\Hsob}[1]{H^{#1}}
\newcommand{\Cinfty}{C^\infty}

% Scalars, vectors, matrices, tensors
\newcommand{\scal}[1]{#1}
\newcommand{\vct}[1]{\symbfit{#1}}   % bold italic vector, works for Greek
\newcommand{\mat}[1]{\symbf{#1}}     % bold upright matrix
\newcommand{\ten}[1]{\symbfsf{#1}}   % bold sans-serif tensor

% Unit vectors
\newcommand{\ihat}{\hat{\vct{i}}}
\newcommand{\jhat}{\hat{\vct{j}}}
\newcommand{\khat}{\hat{\vct{k}}}

% Transpose, inverse, trace, determinant
\newcommand{\T}{^{\mathsf{T}}}
\newcommand{\inv}{^{-1}}
\newcommand{\tr}{\operatorname{tr}}
\newcommand{\diag}{\operatorname{diag}}
\newcommand{\rank}{\operatorname{rank}}
\newcommand{\nullspace}{\operatorname{null}}
\newcommand{\range}{\operatorname{range}}
\newcommand{\determinant}{\operatorname{det}}

% Norms, inner products
\newcommand{\norm}[1]{\left\lVert #1 \right\rVert}
\newcommand{\abs}[1]{\left\lvert #1 \right\rvert}
\newcommand{\inner}[2]{\left\langle #1, #2 \right\rangle}

% Differential operators
\newcommand{\dd}{\,\mathrm{d}}
\newcommand{\pd}{\partial}
\newcommand{\grad}{\nabla}
\newcommand{\divg}{\nabla \cdot}
\newcommand{\curlop}{\nabla \times}
\newcommand{\rot}{\nabla \times}
\newcommand{\laplace}{\Delta}

% Derivatives
\newcommand{\dv}[2]{\frac{\mathrm{d} #1}{\mathrm{d} #2}}
\newcommand{\dvtwo}[2]{\frac{\mathrm{d}^2 #1}{\mathrm{d} #2^2}}
\newcommand{\pdv}[2]{\frac{\partial #1}{\partial #2}}
\newcommand{\pdvtwo}[2]{\frac{\partial^2 #1}{\partial #2^2}}

% Matrix exponential
\newcommand{\expm}{\operatorname{expm}}

% Complex / imaginary numbers
\newcommand{\ii}{\mathrm{i}}
\newcommand{\Real}{\operatorname{Re}}
\newcommand{\Imag}{\operatorname{Im}}

% Fourier transform
\newcommand{\FT}{\mathcal{F}}
\newcommand{\IFT}{\mathcal{F}^{-1}}
\newcommand{\fourier}[1]{\widehat{#1}}

% Probability and statistics
\newcommand{\E}{\mathbb{E}}
\newcommand{\Var}{\operatorname{Var}}
\newcommand{\Cov}{\operatorname{Cov}}
\newcommand{\Prob}{\mathbb{P}}
\newcommand{\Normal}{\mathcal{N}}

% Quaternions
\newcommand{\quat}[1]{\mathbf{#1}}
\newcommand{\qconj}[1]{#1^{*}}
\newcommand{\qinv}[1]{#1^{-1}}
\newcommand{\qmul}{\otimes}

% Stochastic calculus
\newcommand{\Wiener}{W_t}
\newcommand{\ito}{\mathrm{d}}

% Numerical methods
\newcommand{\bigO}{\mathcal{O}}

% ------------------------------------------------------------
% Article-specific macros
% ------------------------------------------------------------
\newcommand{\sat}{\operatorname{sat}}
\newcommand{\wrap}{\operatorname{wrap}_{180}}
\newcommand{\sgn}{\operatorname{sgn}}
\newcommand{\clip}{\operatorname{clip}}
\newcommand{\heading}{\psi}
\newcommand{\cmd}{\mathrm{cmd}}
\newcommand{\xte}{\mathrm{xte}}
\newcommand{\TWA}{\mathrm{TWA}}
\newcommand{\AWA}{\mathrm{AWA}}
\newcommand{\COG}{\mathrm{COG}}
\newcommand{\SOG}{\mathrm{SOG}}

% ------------------------------------------------------------
% Document metadata
% ------------------------------------------------------------
\title{Algorithmic Methods in Open Source Marine Autopilots}
\author{
Mikhail S. Grushinskiy\\
Independent Researcher, USA, New Jersey\\
Bareboat Necessities GitHub Project
}
\date{\today}

\begin{document}

\maketitle
\suppressfloats[t]

\begin{abstract}
This article summarizes the main algorithms found in the uploaded
\texttt{pypilot-master} source tree and compares them with Fenix Autopilot and common
marine and robotics autopilot architectures.  \texttt{pypilot}, written by
Sean D'Epagnier as an open-source marine autopilot, Fenix Autopilot, introduced by
Sergio Pascual as an Arduino-based DIY tiller pilot, and ArduPilot, which grew from
Jordi Mu\~noz and Chris Anderson's late-2000s DIY Drones work into a broad community
autopilot stack, are treated as credited reference designs.  The active
\texttt{pypilot} control path is best understood as a practical marine stack: quaternion attitude estimation and low-pass heading
signals, mode-dependent heading generation, a saturated heading-error controller,
a tack override, and a motor-oriented servo layer.  The default working pilot is not
an advanced nonlinear path follower; it is an explicit extended proportional--derivative
heading controller with gyro-rate, rate-of-rate, square-root heading-error, and
heading-command feed-forward terms.  Compared with ArduPilot/PX4-style stacks, its
path-following layer is much simpler; compared with Fenix Autopilot, it is more
distributed and less tied to one embedded actuator package, but both projects make
motor, rudder, and low-cost drive handling central to the design.
\end{abstract}

\section{Scope and Reading Method}

The analysis is based on the uploaded source archive \texttt{pypilot-master.zip},
with primary attention to the files listed in \cref{tab:source-files}; project designed
by Sean D'Epagnier's \cite{pypilot-author}.
The public pypilot user manual and workbook wiki are used as the main external
pypilot references for operating modes, commissioning, calibration, and user-facing
configuration \cite{pypilot-user-manual,pypilot-workbook-wiki}.  Other external
references are used only to position the design relative to other autopilot families:
ArduPilot Rover navigation and steering-rate control, PX4 fixed-wing position control,
Fenix Autopilot's public repository and documentation, and B\&G descriptions of
commercial sailing autopilot operation.

\begin{table}[!tbp]
\centering
\caption{Important source files inspected in the uploaded tree.}
\label{tab:source-files}
\begin{tabularx}{\columnwidth}{@{}Z{0.34\columnwidth}Y@{}}
\toprule
File & Algorithmic role \\
\midrule
\texttt{autopilot.py} & Main loop, mode selection, heading error, feed-forward command rate. \\
\texttt{pilots/pilot.py} & Base gain accumulation and mode-to-heading conversion. \\
\texttt{pilots/basic.py} & Default extended heading controller. \\
\texttt{pilots/wind.py} & Wind-relative heading generation and wind-gust term. \\
\texttt{pilots/absolute.py} & Rudder-position command experiment. \\
\texttt{tacking.py} & Tack/jibe state machine and override command. \\
\texttt{boatimu.py} & RTIMU-based quaternion attitude, heading/rate filters, alignment. \\
\texttt{gps\_filter.py} & Optional GPS--inertial Kalman filter for position, speed, track. \\
\texttt{sensors.py} & APB/NAV cross-track-to-heading command. \\
\texttt{servo.py} & Motor command shaping, rudder feedback, current/fault protection. \\
\bottomrule
\end{tabularx}
\end{table}

\section{High-Level Architecture}

The control stack has five layers:
\begin{enumerate}
\item inertial and heading estimation;
\item mode-specific definition of the controlled heading variable;
\item heading-command and heading-error calculation;
\item pilot-specific control law;
\item servo/motor command shaping and protection.
\end{enumerate}

A useful abstraction is shown in \cref{fig:pypilot-stack}.  The important point is that
the servo layer is not merely an actuator output.  In this code base it is a significant
part of the closed-loop behavior because the usual command is closer to a rudder-drive
speed request than to an absolute rudder angle.

\begin{figure*}[t]
\centering
\begin{adjustbox}{max width=\textwidth}
\begin{tikzpicture}[
  node distance=0.45cm and 0.65cm,
  block/.style={draw, rounded corners, align=center, minimum height=0.82cm, text width=2.75cm},
  wide/.style={draw, rounded corners, align=center, minimum height=0.82cm, text width=3.2cm},
  arrow/.style={-Latex, thick}
]
\node[block] (imu) {IMU / RTIMU\\quaternion attitude};
\node[block, right=of imu] (signals) {Heading, rate,\\rate-of-rate filters};
\node[block, right=of signals] (mode) {Mode mapping\\compass/GPS/NAV\\wind};
\node[block, right=of mode] (error) {Heading command\\and saturated error};
\node[block, right=of error] (pilot) {Pilot law\\basic/wind/absolute};
\node[wide, right=of pilot] (servo) {Servo layer\\pulse, current, fault, rudder};
\node[block, right=of servo] (boat) {Boat, drive,\\rudder, waves};

\draw[arrow] (imu) -- (signals);
\draw[arrow] (signals) -- (mode);
\draw[arrow] (mode) -- (error);
\draw[arrow] (error) -- (pilot);
\draw[arrow] (pilot) -- (servo);
\draw[arrow] (servo) -- (boat);
\draw[arrow]
  ([yshift=-0.06cm]boat.south)
  to[out=-65,in=-115,looseness=0.55]
  node[below, midway, font=\scriptsize] {motion feedback}
  ([yshift=-0.06cm]imu.south);

\node[block, below=1.0cm of error] (tack) {Tack/jibe state\\machine override};
\draw[arrow] (tack) -- (pilot.south);

\node[block, above=0.9cm of mode] (nmea) {GPS, APB, wind,\\water speed};
\draw[arrow] (nmea) -- (mode);
\end{tikzpicture}
\end{adjustbox}
\caption{Conceptual \texttt{pypilot} control stack in the uploaded source.}
\label{fig:pypilot-stack}
\end{figure*}

\section{Pilot Loading and Active Algorithms}

Pilot modules are discovered dynamically from \texttt{pypilot/pilots}.  A module is
skipped when it declares \texttt{disabled = True}; otherwise a module-level symbol
\texttt{pilot} is appended to the list of available pilot classes.  In the uploaded tree,
the practical active pilots are \texttt{basic}, \texttt{wind}, and \texttt{absolute}; most
adaptive or experimental pilots are disabled or incomplete.  The default selected pilot
is \texttt{basic}.

\begin{table}[!tbp]
\centering
\caption{Pilot modules and operational status in the inspected archive.}
\label{tab:pilots}
\begin{tabularx}{\columnwidth}{@{}Z{0.19\columnwidth}cZ{0.24\columnwidth}Y@{}}
\toprule
Pilot & Active & Class & Main idea \\
\midrule
\texttt{basic} & yes & \texttt{BasicPilot} & Extended heading PD with feed-forward. \\
\texttt{wind} & yes & \texttt{WindPilot} & Wind-relative heading plus gust term. \\
\texttt{absolute} & yes & \texttt{AbsolutePilot} & Rudder-position command, requires rudder feedback. \\
\texttt{simple} & no & \texttt{SimplePilot} & Classic PID, disabled. \\
\texttt{rate} & no & \texttt{RatePilot} & Turn-rate target loop, disabled. \\
\texttt{gps} & no & \texttt{GPSPilot} & GPS-heading pilot, disabled. \\
\texttt{deadzone} & no & \texttt{DeadZonePilot} & Deadband experiment, disabled. \\
\texttt{fuzzy} & no & \texttt{FuzzyPilot} & Online lookup-table learning, disabled. \\
\texttt{vmg} & no & \texttt{VMGPilot} & VMG table idea, incomplete and disabled. \\
\texttt{learning} & no & \texttt{LearningPilot} & TFLite predictive idea, disabled. \\
\bottomrule
\end{tabularx}
\end{table}

This distinction matters.  The repository contains sketches of learning and VMG control,
but the real steering path in this archive remains deterministic and comparatively simple.

\section{Heading Representation and Error Calculation}

Let \(\heading\) denote the current controlled heading variable and
\(\heading_{\cmd}\) the commanded value.  Depending on mode, \(\heading\) is not always
magnetic compass heading.  The base pilot maps sensor data into a mode-dependent heading:
\begin{align}
\heading_{\mathrm{compass}} &= \heading_c, \\
\heading_{\mathrm{gps}} &= \wrap(\heading_c + \Delta_{g/c}), \\
\heading_{\mathrm{wind}} &= \wrap(\Delta_{w/c} - \heading_c), \\
\heading_{\mathrm{true\ wind}} &= \wrap(\Delta_{tw/c} - \heading_c),
\end{align}
where \(\Delta_{g/c}\), \(\Delta_{w/c}\), and \(\Delta_{tw/c}\) are slowly updated offsets
between compass, GPS course, apparent wind, and true wind references.

The signed angular wrapping operation is represented by
\begin{equation}
  \wrap(a; b) \in (b-180^{\circ}, b+180^{\circ}],
\end{equation}
matching the source function \texttt{resolv(angle, offset)}.  The heading error is then
clipped:
\begin{equation}
  e = \clip\!\left(\wrap(\heading - \heading_{\cmd}), -30^{\circ}, 30^{\circ}\right).
  \label{eq:error-clip}
\end{equation}
For wind modes the sign is reversed, because apparent wind direction is defined as where
the wind is coming from, not where the vessel is traveling.

This saturation is a practical marine choice.  It prevents the inner controller from
seeing very large error after a sensor jump, tack, gybe, compass recalibration, or command
change.  It also means \texttt{pypilot} prefers bounded correction over aggressively
trying to unwind a large angular error.

\section{The Default Basic Pilot}

The default basic pilot computes heading rate \(\dot{\heading}\) and heading-rate-rate
\(\ddot{\heading}\) from the IMU-derived signals
\texttt{headingrate\_lowpass} and \texttt{headingraterate\_lowpass}.  The active control
law can be written as
\begin{equation}
\begin{split}
  u_b ={}& K_P e
        + K_D \dot{\heading}
        + K_{DD} \ddot{\heading}
        + K_{PR}\,\sgn(e)\sqrt{\abs{e}} \\
        &+ K_{FF}\dot{\heading}_{\cmd} .
\end{split}
\label{eq:basic}
\end{equation}

Here \(u_b\) is the requested servo command before the lower servo layer reshapes it.  The
square-root term is computed explicitly from heading error:
\begin{equation}
  e_{PR} = \sgn(e)\sqrt{\abs{e}}.
  \label{eq:pr}
\end{equation}

The terms in \cref{eq:basic} have intuitive roles:
\begin{itemize}
\item \(K_P e\): proportional correction from course error;
\item \(K_D\dot{\heading}\): gyro counter-rudder or yaw-rate damping;
\item \(K_{DD}\ddot{\heading}\): damping against rapidly changing yaw rate;
\item \(K_{PR}e_{PR}\): nonlinear low-error authority without a purely high linear gain;
\item \(K_{FF}\dot{\heading}_{\cmd}\): feed-forward during commanded course changes.
\end{itemize}

The interesting part is the combination of \(D\), \(DD\), and \(PR\).  A commercial pilot
might expose related concepts as \emph{gain}, \emph{counter rudder}, \emph{response}, or
\emph{sea state}; \texttt{pypilot} exposes the actual additive terms.

\section{Integral Action and Bias Handling}

The main autopilot maintains an integrated heading error,
\begin{equation}
  I_{k+1} = \clip\!\left(I_k + \frac{e_k}{50}\Delta t, -10, 10\right),
  \label{eq:integral}
\end{equation}
but the default basic and wind pilots comment out the \(I\) term.  The code also reduces
or resets this integral around mode changes and command changes.  This is conservative:
it avoids integral windup in waves and maneuvers, but persistent biases such as weather
helm, current, or asymmetric drive behavior are not estimated as an explicit disturbance
state in the active pilot.

A more model-based heading controller would estimate a slowly varying yaw bias,
\begin{equation}
  b_{k+1} = b_k + \eta_k,
\end{equation}
and subtract it from the required control.  In this archive, comparable behavior is
handled indirectly through offsets, manual gain tuning, servo calibration, and the lower
servo windup/pulse mechanism.

\section{Heading-Command Feed-Forward}

When the autopilot is enabled, the main loop computes a filtered heading-command-rate
signal from the difference between the current and previous commands:
\begin{equation}
\begin{split}
  \dot{\heading}_{\cmd,k} ={}&
  (1-\alpha)\dot{\heading}_{\cmd,k-1} \\
  &+ \alpha\,\wrap(\heading_{\cmd,k}-\heading_{\cmd,k-1}), \\
  &\alpha \approx 0.1.
\end{split}
\end{equation}
For non-wind modes the sign is reversed before filtering.  This gives the controller a
feed-forward term during user-initiated course changes or tack state transitions.  It is
not a reference-model trajectory generator; it is a first-order filtered command-difference
signal.

\section{Wind Pilot}

The wind pilot preserves the same general heading-control structure, but it changes how
the controlled heading is defined.  In apparent wind mode,
\begin{equation}
  \heading = \AWA,
\end{equation}
while in GPS or compass modes it can synthesize a GPS- or compass-like heading from wind
offsets.  It also computes a wind-gust signal
\begin{equation}
  g_w(k) = V_w(k)-V_w(k-1),
\end{equation}
with a sign adjustment depending on wind direction.  The wind-pilot command is therefore
approximately
\begin{equation}
  u_w = K_P e + K_D\dot{\heading} + K_{DD}\ddot{\heading} + K_{WG}g_w.
  \label{eq:wind}
\end{equation}

This is a useful sailing-specific heuristic, but it is not a full aerodynamic model.  It
does not estimate sail force, rudder lift, polar performance, or wave phase.  It is a
simple, inspectable wind-aware heading controller.

\section{NAV/APB Mode}

The APB sensor handler implements a very compact route-following law.  If an APB sentence
provides route track \(\chi_r\) and cross-track error \(x_{te}\), then the heading command
is
\begin{equation}
  \heading_{\cmd} = \chi_r + K_{xte} x_{te}.
  \label{eq:apb}
\end{equation}
The source default is \(K_{xte}=300\), documented in a comment as roughly
\(30^{\circ}\) for \(0.1\) nautical mile.

This design is much simpler than modern path-following controllers.  It has no explicit
look-ahead distance, no velocity-dependent lateral-acceleration target, no curvature
planning, and no inner yaw-rate target.  The advantage is simplicity and compatibility
with common NMEA/APB navigation data.  The disadvantage is weaker behavior in current,
tight routes, variable speed, and overshoot-prone waypoint transitions.

\section{Tack and Gybe State Machine}

The tack module is a supervisory override.  Its states are
\begin{equation}
  \mathcal{S}=\{\texttt{none},\texttt{begin},\texttt{waiting},\texttt{tacking}\}.
\end{equation}
The tack direction is inferred from wind direction when possible, or from heel if enabled.
During \texttt{tacking}, the normal pilot is bypassed and the servo receives
\begin{equation}
  u_t = \clip\!\left(1 + \frac{\dot{\heading}+\ddot{\heading}/2}{r_t}, -1, 1\right)
\end{equation}
up to a direction sign, where \(r_t\) is the configured tack rate.  Completion is detected
when a normalized progress variable exceeds a configured threshold.  The command heading
is then switched to the new tack heading.

This is intentionally simple: it is a maneuver state machine, not an optimal tack planner.
It does not time the tack using waves or boat speed, and it does not evaluate alternative
tacking strategies.  Adaptive tack-strategy selection is a separate research topic in
robotic sailing literature and is not implemented in the active \texttt{pypilot} code path.

\section{Servo and Motor Command Layer}

The servo layer is one of the most distinctive parts of the project.  A pilot output
\(u\) is not simply sent to a rudder.  The normal command path treats \(u\) as a desired
motor-drive speed-like request, then reshapes it using minimum speed, maximum speed,
periodic pulse behavior, current limits, voltage compensation, overcurrent end-stop
protection, and rudder feedback when available.

Let \(v_c\) be the desired drive speed after pilot gain.  The layer applies
\begin{equation}
  v_c \leftarrow \clip(v_c, -v_{\max}, v_{\max}),
\end{equation}
and a minimum-motion pulse mechanism.  The source maintains a windup-like accumulator
\begin{equation}
  W_{k+1} = W_k + (v_c - v_k)\Delta t,
\end{equation}
and only emits a nonzero command when the accumulated request is large enough to overcome
the configured minimum drive speed over the configured period.  This is a practical answer
to real marine actuators: friction, deadband, clutch delay, slow hydraulic response, and
cheap H-bridge motor control.

When rudder feedback is available, the \texttt{absolute} pilot and servo position command
path can use a rudder-position PID-like loop:
\begin{equation}
  u_r = K_{rp}(r_{\cmd}-r) + K_{ri}\bar e_r + K_{rd}\dot r.
\end{equation}
However, the default \texttt{basic} pilot does not require rudder feedback.

\section{IMU and Heading Estimation}

The IMU layer uses RTIMU quaternion fusion, applies alignment calibration, and computes
Euler roll, pitch, and heading from the aligned quaternion.  Quaternion rotation can be
written as
\begin{equation}
  \vct{v}' = \quat{q}\qmul \vct{v}\qmul \qconj{\quat{q}}.
\end{equation}
The source then derives heading rate from the rotated gyro vector and computes low-pass
signals:
\begin{align}
  \heading_{lp,k} &= \alpha_h \heading_k + (1-\alpha_h)\heading_{lp,k-1}, \\
  \dot{\heading}_{lp,k} &= \alpha_r \dot{\heading}_k
     + (1-\alpha_r)\dot{\heading}_{lp,k-1}, \\
  \ddot{\heading}_{lp,k} &= \alpha_a \ddot{\heading}_k
     + (1-\alpha_a)\ddot{\heading}_{lp,k-1}.
\end{align}
Angular wrap is handled specially for heading low-pass filtering so that filtering across
\(0^{\circ}/360^{\circ}\) does not introduce a false large jump.

The heading estimator is not a full vessel INS.  It is primarily an AHRS with marine
alignment and calibration behavior, feeding the heading-control loop.

\section{Optional GPS--Inertial Filter}

The optional GPS filter estimates position and velocity at a higher output rate by
combining GPS updates with acceleration-derived prediction.  In two dimensions the state
can be written as
\begin{equation}
  \vct{x} = \begin{bmatrix} x & y & v_x & v_y \end{bmatrix}\T .
\end{equation}
The prediction model is the constant-acceleration kinematic model
\begin{equation}
  \vct{x}_{k+1} =
  \begin{bmatrix}
  \mat{I} & \Delta t\mat{I} \\
  \mat{0} & \mat{I}
  \end{bmatrix}\vct{x}_k
  +
  \begin{bmatrix}
  \frac{1}{2}\Delta t^2\mat{I} \\
  \Delta t\mat{I}
  \end{bmatrix}\vct{a}_k.
  \label{eq:gps-predict}
\end{equation}
The GPS update uses a standard linear Kalman correction
\begin{align}
  \vct{y}_k &= \vct{z}_k - \mat{H}\vct{x}_k, \\
  \mat{S}_k &= \mat{H}\mat{P}_k\mat{H}\T + \mat{R}, \\
  \mat{K}_k &= \mat{P}_k\mat{H}\T\mat{S}_k\inv, \\
  \vct{x}_k^+ &= \vct{x}_k^- + \mat{K}_k\vct{y}_k, \\
  \mat{P}_k^+ &= (\mat{I}-\mat{K}_k\mat{H})\mat{P}_k^-.
\end{align}
The implementation also keeps a history so that delayed GPS measurements can be applied
at the approximate measurement time and then fast-forwarded through later inertial
predictions.

This is useful, but it is not the core steering controller.  It mostly improves filtered
speed, track, and position availability.

\section{Comparison with ArduPilot Rover}

ArduPilot began in the DIY Drones community around Jordi Mu\~noz and Chris Anderson's
Arduino-era autopilot work and later evolved, through the ArduPilot development team and
contributors such as Andrew ``Tridge'' Tridgell, Pat Hickey, Randy Mackay, and many
others, into a multi-vehicle open-source autopilot project \cite{ardupilot-history}.
ArduPilot Rover separates high-level path following from low-level steering-rate control.
Its older navigation documentation describes the L1 controller as the highest-level
steering controller: it accepts position targets and current speed and outputs a desired
lateral acceleration, which then feeds lower-level controllers \cite{ardupilot-l1}.  Its
steering-rate documentation describes a turn-rate PID controller with feed-forward,
proportional, integral, and derivative gains \cite{ardupilot-rate}.

A compact comparison is
\begin{align}
\texttt{pypilot:}\quad
& \chi_r + K_{xte}x_{te} \to \heading_{\cmd} \to u_b, \\
\text{ArduPilot-style:}\quad
& \text{path} \to a_y^{\cmd}, \\
& a_y^{\cmd} \to \dot{\heading}^{\cmd} \to u.
\end{align}

Thus \texttt{pypilot} is heading-hold centered, whereas ArduPilot Rover is path-following
centered in autonomous modes.  The \texttt{pypilot} approach is lighter and easier to tune
on a sailboat with limited sensors, but it is less rigorous for route tracking.

\section{Comparison with PX4}

PX4 fixed-wing control is even more layered.  PX4 documentation describes TECS for
altitude and airspeed control and NPFG for horizontal heading/position control
\cite{px4-fw}.  That architecture is built for aircraft energy and trajectory control,
with separate inner attitude and rate loops.  A direct comparison to \texttt{pypilot} is
therefore imperfect, but the contrast is useful:
\begin{itemize}
\item PX4 treats the vehicle as a dynamic aircraft with explicit position, attitude,
energy, and actuator layers.
\item \texttt{pypilot} treats the sailboat autopilot mainly as a heading servo plus
marine-specific modes and motor protection.
\end{itemize}

The difference is not just sophistication.  The vehicles are different.  A small sailboat
autopilot often needs low power, reliable clutch/motor handling, wind modes, and tolerance
of slow response more than it needs aggressive trajectory tracking.

\section{Comparison with Commercial Marine Pilots}

The pypilot user manual and workbook present \texttt{pypilot} as a modular
open-source marine autopilot supporting compass, GPS/NAV, wind, and true-wind
operation, with IMU calibration, optional GPS/wind inputs, network clients, and
motor-controller commissioning \cite{pypilot-user-manual,pypilot-workbook-wiki}.
B\&G pilot materials similarly describe wind and navigation modes and manual course
adjustments from pilot controls \cite{bg-pilot}.  B\&G sailing guidance also
describes true-wind-angle use for gybing \cite{bg-setup}.

The practical difference is visibility.  Commercial pilots may provide polished
commissioning, proprietary sea-state adaptation, and integrated displays, but the actual
control laws are hidden.  \texttt{pypilot} exposes the terms, gains, and source code.
That is a major advantage for experimentation, but it also means incomplete experimental
code is visible in the tree.

\section{Experimental and Incomplete Algorithms}

The repository contains several ideas that are not part of the active default control path.

\subsection{VMG Table Pilot}

The VMG pilot attempts to collect measurements of speed and track versus heading on each
tack, storing them in heading-indexed tables.  Conceptually, one would estimate
\begin{equation}
  \mathrm{VMG}(\heading) = \vct{v}(\heading)\cdot \hat{\vct{d}},
\end{equation}
where \(\hat{\vct{d}}\) is the desired progress direction.  The command would then move
toward
\begin{equation}
  \heading^* = \arg\max_{\heading} \mathrm{VMG}(\heading).
\end{equation}
In the uploaded file this optimizer is not implemented; the relevant block contains a
\texttt{pass}.  The module is also disabled.

\subsection{Fuzzy Pilot}

The fuzzy pilot defines a multi-dimensional interpolation table over ground speed, wind
speed, wind direction, rudder angle, heel, sea state, heading error, and heading rate.
It recursively interpolates the table to compute a command and updates old table entries
from a delayed error signal.  The learning update is approximately
\begin{equation}
  c_{new} = \clip(c_{old} + \lambda\epsilon, -1, 1),
\end{equation}
where \(\epsilon\) is a delayed heading/rate error proxy.  This pilot is disabled.

\subsection{Learning Pilot}

The learning pilot sketches a TensorFlow Lite model-predictive idea: feed a history of
sensor values and candidate future actions to a learned model, predict future outputs, and
choose the action sequence with lowest loss,
\begin{equation}
  J = \sum_i
  \left(K_P \hat e_i + K_D\widehat{\dot{\heading}}_i\right)^2
  + K_W u_i^2.
\end{equation}
This is promising in principle, but in the uploaded tree it is disabled and incomplete.

\section{Code-Level Observations}

The following issues are important if these experimental paths are to be revived:
\begin{itemize}
\item \texttt{pilots/vmg.py} contains a misspelling, \texttt{heading\_comand}, where
\texttt{heading\_command} is intended.
\item \texttt{pilots/vmg.py} leaves the core optimum-direction calculation unimplemented.
\item \texttt{pilots/autotune.py} references \texttt{cost} in places where
\texttt{self.cost} appears intended.
\item \texttt{pilots/gps.py} calls \texttt{resolve(...)} although the resolver function
in the uploaded tree is named \texttt{resolv(...)}.
\item The active \texttt{wind.py} true-wind branch references \texttt{ap.true\_wind\_sensor};
this property was not found elsewhere in the inspected source tree.
\end{itemize}

These observations do not invalidate the active \texttt{basic} pilot, but they show that
the repository contains research sketches as well as production code.

\section{Strengths of the Design}

The main strengths are:
\begin{enumerate}
\item \textbf{Explicit control law.}  The default controller is mathematically simple and
visible in the source.
\item \textbf{Excellent practical actuator attention.}  The servo layer handles motor and
rudder realities that many abstract autopilot algorithms ignore.
\item \textbf{Marine-specific heading modes.}  Compass, GPS/NAV, wind, and true-wind
concepts are built into the architecture.
\item \textbf{Low computational burden.}  The core controller is suitable for Raspberry Pi
class hardware and modest update rates.
\item \textbf{Hackability.}  Gains, modes, and experimental pilots can be inspected and
changed without reverse engineering a proprietary pilot.
\end{enumerate}

\section{Weaknesses and Upgrade Direction}

The main weakness is not the basic heading loop; it is the lack of a stronger guidance
layer above it.  NAV mode should be upgraded before any machine-learning approach is
attempted.  A natural deterministic roadmap is:
\begin{enumerate}
\item keep the existing servo protection and calibration layer;
\item keep the AHRS and heading-mode architecture;
\item replace \cref{eq:apb} with an L1, pure-pursuit, or Stanley-style path follower;
\item add a speed-aware yaw-rate target layer;
\item use rudder-angle feedback when available;
\item add a slow weather-helm/current bias estimator;
\item only then revive VMG/polar optimization.
\end{enumerate}

A possible marine guidance layer would compute a desired lateral acceleration
\begin{equation}
  a_y^{\cmd} = f(\text{path}, \vct{p}, \vct{v}, L_1),
\end{equation}
convert it to a desired turn rate
\begin{equation}
  \dot{\heading}^{\cmd} \approx \frac{a_y^{\cmd}}{\max(V, V_{min})},
\end{equation}
and then use a rate loop or extended heading loop to command the existing servo layer.
This keeps the best part of \texttt{pypilot}---its actuator realism---while improving
route tracking.

\section{Low-Level Servo Control Protocol}
\label{sec:servo-protocol}
The pypilot software does not command the motor controller through a
textual marine protocol such as NMEA~0183 or Signal~K.  The normal
Arduino motor-controller path uses a compact binary serial protocol.
This protocol is implemented by the host-side Arduino servo driver and
the microcontroller firmware.  It is deliberately small: every message is
a fixed-length packet carrying one command code, one 16-bit value, and a
CRC byte.
\subsection{Transport and Packet Format}
The physical transport is UART serial.  The pypilot host probes the
Arduino servo controller at approximately \(38400\) baud.  The Arduino
firmware uses the usual serial byte framing, effectively \(8\)-data-bit,
no-parity, one-stop-bit framing.  The application protocol itself is not
stream-delimited by a start byte.  Instead, synchronization is obtained
statistically by requiring valid CRC-protected packets.
Each packet has exactly four bytes:
\[
\begin{array}{c|c|c|c}
b_0 & b_1 & b_2 & b_3 \\
\hline
\text{code} & \text{value}_{\mathrm{lo}} &
\text{value}_{\mathrm{hi}} & \operatorname{crc8}(b_0,b_1,b_2)
\end{array}.
\]
The 16-bit payload is little-endian:
\[
v = b_1 + 256 b_2.
\]
The CRC is computed over the first three bytes only:
\[
c_0 = \mathtt{0xff},
\qquad
c_{k+1} =
\operatorname{table}\!\left[c_k \oplus b_k\right],
\qquad
b_3 = c_3 .
\]
The table used in the source corresponds to a CRC-8 polynomial
\(\mathtt{0x31}\) with initial value \(\mathtt{0xff}\).  There is no
length field, sequence number, escaping, or acknowledgement byte.  On bad
CRC, the firmware clears synchronization and stops the motor.  The
controller requires multiple valid packets before accepting the stream as
synchronized.
\subsection{Main Motor Command}
The most important host-to-controller packet is
\texttt{COMMAND\_CODE}, whose command byte is
\[
b_0 = \mathtt{0xC7}.
\]
The pypilot command is a normalized scalar
\[
u \in [-1,1],
\]
where the sign selects the motor direction and the magnitude selects the
drive effort.  The value placed into the packet is
\[
v = 1000(u+1).
\]
Therefore,
\[
\begin{array}{lll}
u=-1 &\Rightarrow& v=0,\\
u=0 &\Rightarrow& v=1000,\\
u=1 &\Rightarrow& v=2000.
\end{array}
\]
The neutral command is thus not zero, but \(v=1000\).  The firmware
rejects command values above \(2000\).  A compatible host can therefore
send a motor command by constructing
\[
\begin{aligned}
b_0 &= \mathtt{0xC7},\\
b_1 &= v \bmod 256,\\
b_2 &= \left\lfloor v/256 \right\rfloor,\\
b_3 &= \operatorname{crc8}(b_0,b_1,b_2).
\end{aligned}
\]
For example, the neutral command has
\[
v = 1000 = \mathtt{0x03E8},
\]
so the first three bytes are
\[
\mathtt{0xC7},\quad \mathtt{0xE8},\quad \mathtt{0x03}.
\]
\subsection{Host-to-Controller Commands}
\begin{table*}[!tbp]
\centering
\caption{Host-to-servo packet codes used by the pypilot Arduino motor-controller protocol.}
\label{tab:pypilot-servo-host-codes}
\begin{tabularx}{\textwidth}{Z{0.24\textwidth}Z{0.10\textwidth}Z{0.20\textwidth}Y}
\toprule
Name & Code & Payload & Purpose \\
\midrule
\texttt{COMMAND\_CODE} &
\(\mathtt{0xC7}\) &
\(0\ldots 2000\) &
Main motor command, with \(1000\) as neutral. \\
\texttt{RESET\_CODE} &
\(\mathtt{0xE7}\) &
ignored &
Clear overcurrent fault state. \\
\texttt{MAX\_CURRENT\_CODE} &
\(\mathtt{0x1E}\) &
current limit &
Set motor current limit, approximately in \(10\,\mathrm{mA}\) units. \\
\texttt{MAX\_CONTROLLER\_TEMP\_CODE} &
\(\mathtt{0xA4}\) &
\(100T\) &
Set controller over-temperature threshold in centi-degrees Celsius. \\
\texttt{MAX\_MOTOR\_TEMP\_CODE} &
\(\mathtt{0x5A}\) &
\(100T\) &
Set motor over-temperature threshold in centi-degrees Celsius. \\
\texttt{RUDDER\_MIN\_CODE} &
\(\mathtt{0x2B}\) &
encoded rudder position &
Set the minimum rudder stop. \\
\texttt{RUDDER\_MAX\_CODE} &
\(\mathtt{0x4D}\) &
encoded rudder position &
Set the maximum rudder stop. \\
\texttt{REPROGRAM\_CODE} &
\(\mathtt{0x19}\) &
ignored &
Enter bootloader/reprogramming mode. \\
\texttt{DISENGAGE\_CODE} &
\(\mathtt{0x68}\) &
ignored &
Stop and disengage the motor output. \\
\texttt{MAX\_SLEW\_CODE} &
\(\mathtt{0x71}\) &
two 8-bit fields &
Set fast and slow slew-rate limits. \\
\texttt{EEPROM\_READ\_CODE} &
\(\mathtt{0x91}\) &
start/end address &
Request EEPROM bytes from the controller. \\
\texttt{EEPROM\_WRITE\_CODE} &
\(\mathtt{0x53}\) &
address/value byte &
Write one byte to EEPROM. \\
\texttt{CLUTCH\_PWM\_AND\_BRAKE\_CODE} &
\(\mathtt{0x36}\) &
PWM/brake byte pair &
Set clutch PWM and brake-enable flag. \\
\bottomrule
\end{tabularx}
\end{table*}
The \texttt{DISENGAGE\_CODE} is important because the protocol is
designed around continuous command streaming.  A command is not treated as
a permanent latched actuator state.  If the command stream becomes invalid
or stale, the controller is expected to stop or disengage the motor.
\subsection{Rudder Encoding}
Rudder feedback and rudder limits use a normalized position coordinate
approximately in the interval
\[
r \in [-0.5,0.5].
\]
When pypilot sends a rudder stop to the controller, it encodes the
position as
\[
v = 65472(r+0.5).
\]
The inverse decoding used for telemetry is
\[
r = \frac{v}{65472} - 0.5.
\]
The special value
\[
v = 65535
\]
indicates invalid or unavailable rudder feedback.
\subsection{Controller-to-Host Telemetry}
The motor controller returns telemetry using the same four-byte packet
format.  The first byte is a status code, the next two bytes form a
little-endian value, and the fourth byte is the CRC.
\begin{table}[!t]
\centering
\caption{Servo telemetry packet codes.}
\label{tab:pypilot-servo-telemetry}
\begin{tabularx}{\columnwidth}{Z{0.42\columnwidth}Z{0.17\columnwidth}Y}
\toprule
Name & Code & Decoding \\
\midrule
\texttt{CURRENT\_CODE} &
\(\mathtt{0x1C}\) &
\(I = v/100\) amperes. \\
\texttt{VOLTAGE\_CODE} &
\(\mathtt{0xB3}\) &
\(V = v/100\) volts. \\
\texttt{CONTROLLER\_TEMP\_CODE} &
\(\mathtt{0xF9}\) &
Signed \(v/100\) degrees Celsius. \\
\texttt{MOTOR\_TEMP\_CODE} &
\(\mathtt{0x48}\) &
Signed \(v/100\) degrees Celsius. \\
\texttt{RUDDER\_SENSE\_CODE} &
\(\mathtt{0xA7}\) &
\(r = v/65472 - 0.5\), except \(65535\) means invalid. \\
\texttt{FLAGS\_CODE} &
\(\mathtt{0x8F}\) &
Raw 16-bit controller flag word. \\
\texttt{EEPROM\_VALUE\_CODE} &
\(\mathtt{0x9A}\) &
EEPROM address/value byte pair. \\
\bottomrule
\end{tabularx}
\end{table}
The telemetry schedule is asymmetric: current and rudder position are
reported more often than temperature.  Flags are also emitted regularly
so that the host can detect synchronization, engagement, faults, and
controller reboot state.
\subsection{Controller Flags}
The \texttt{FLAGS\_CODE} packet carries a 16-bit hardware-controller
status word.  The important bits are listed in
\cref{tab:pypilot-servo-flags}.
\begin{table}[!t]
\centering
\caption{Servo-controller hardware flags.}
\label{tab:pypilot-servo-flags}
\begin{tabularx}{\columnwidth}{Z{0.16\columnwidth}Y}
\toprule
Value & Meaning \\
\midrule
\(1\) & Serial stream synchronized. \\
\(2\) & Over-temperature fault. \\
\(4\) & Over-current fault. \\
\(8\) & Motor output engaged. \\
\(16\) & Invalid packet or CRC error observed. \\
\(32\) & Port-side output pin fault. \\
\(64\) & Starboard-side output pin fault. \\
\(128\) & Bad supply-voltage fault. \\
\(256\) & Minimum rudder limit reached. \\
\(512\) & Maximum rudder limit reached. \\
\(1024\) & Current-range indicator. \\
\(2048\) & Fuse fault. \\
\(32768\) & Controller rebooted. \\
\bottomrule
\end{tabularx}
\end{table}
Additional flags used inside pypilot, such as driver timeout or saturated
output, are host-side software flags.  They are not transmitted as part of
the 16-bit Arduino status word.
\subsection{Design Consequences}
The protocol is intentionally minimal, but it strongly shapes the
autopilot architecture.  The higher-level pilot computes a normalized
actuator command, while the servo layer is responsible for protecting the
drive electronics and motor.  This separation is important: pypilot does
not assume an ideal rudder actuator.  Instead, it treats the motor
controller as a guarded device with current limits, voltage checks,
temperature limits, rudder stops, end-stop detection, and synchronization
loss behavior.
This makes the pypilot servo protocol closer to a robust embedded motor
link than to a general navigation protocol.  It carries no waypoints,
routes, wind data, or heading commands.  It carries only low-level actuator
commands, configuration values, and motor-controller telemetry.  The
result is a simple binary link that is easy to implement on a small
microcontroller while still being safe enough for a continuously running
marine autopilot actuator.

\section{TCP Command and Value Protocol}
\label{sec:pypilot-tcp-protocol}
In addition to the low-level Arduino servo protocol described in
\cref{sec:servo-protocol}, pypilot exposes a higher-level TCP protocol for
autopilot control, monitoring, configuration, and client synchronization.
This TCP protocol is the normal interface used by pypilot clients, web
controls, remote buttons, display panels, and external software that wants
to command the autopilot.  It is not a binary actuator protocol.  Instead,
it is a newline-delimited textual value protocol.
The low-level servo protocol carries motor-drive commands such as
``drive port'', ``drive starboard'', current limits, voltage, and rudder
feedback.  The TCP protocol carries semantic autopilot values such as
\texttt{ap.enabled}, \texttt{ap.mode}, \texttt{ap.heading\_command},
\texttt{ap.heading}, and \texttt{ap.pilot}.  In other words, the TCP
protocol talks to the autopilot state machine, while the servo protocol
talks to the motor controller.
\subsection{Transport}
The pypilot server opens a TCP listener on all interfaces:
\[
\texttt{0.0.0.0:23322}.
\]
The default port is
\[
p_{\mathrm{tcp}} = 23322,
\]
although the server source allows the port to be overridden with the
environment variable \texttt{PYPILOT\_PORT}.  The normal pypilot clients
use port \texttt{23322}.  The server also advertises itself by Zeroconf
using the service type
\[
\texttt{\_pypilot.\_tcp.local.}
\]
so that clients can discover the active pypilot host on a local network.
The transport characteristics are summarized in
\cref{tab:pypilot-tcp-transport}.
\begin{table}[!t]
\centering
\caption{pypilot TCP value protocol transport.}
\label{tab:pypilot-tcp-transport}
\begin{tabularx}{\columnwidth}{Z{0.36\columnwidth}Y}
\toprule
Field & Value \\
\midrule
Transport & TCP stream \\
Default port & \texttt{23322} \\
Bind address & \texttt{0.0.0.0} \\
Discovery & Zeroconf, \texttt{\_pypilot.\_tcp.local.} \\
Encoding & UTF-8/ASCII-compatible text \\
Framing & newline-delimited records \\
Record form & \texttt{name=json-value\textbackslash n} \\
Security & no authentication or TLS in the protocol itself \\
\bottomrule
\end{tabularx}
\end{table}
Because the protocol has no authentication or encryption layer, it should
be treated as a trusted-LAN interface.  Any host able to connect to the
TCP port can attempt to modify writable autopilot values.
\subsection{Line Grammar}
Each protocol record is a single assignment line:
\[
\texttt{name=json-value\textbackslash n}.
\]
The left side is a pypilot value name.  The right side is a JSON value.
The whole line is not JSON; only the right-hand side is JSON.
A compact grammar is
{\scriptsize
\[
\begin{aligned}
\text{record} &::= \text{name}\; \texttt{=}\; \text{json-value}\; \texttt{\textbackslash n},\\
\text{name} &::= \text{dotted pypilot value name},\\
\text{json-value} &::= \text{number} \mid \text{string} \mid \text{boolean}
                   \mid \text{array} \mid \text{object}.
\end{aligned}
\]
}
Thus a valid command line is
\[
\texttt{ap.enabled=true\textbackslash n},
\]
not
\[
\texttt{\{"ap.enabled":true\}\textbackslash n}.
\]
The right-hand side follows JSON conventions.  Booleans are
\texttt{true} and \texttt{false}, strings must be quoted, and objects use
normal JSON syntax.  Several examples are shown in
\cref{tab:pypilot-tcp-examples}.
\begin{table}[!t]
\centering
\caption{Examples of valid pypilot TCP protocol records.}
\label{tab:pypilot-tcp-examples}
\begin{tabularx}{\columnwidth}{Z{0.36\columnwidth}Y}
\toprule
Purpose & Record \\
\midrule
Engage autopilot &
\texttt{ap.enabled=true} \\
Disengage autopilot &
\texttt{ap.enabled=false} \\
Set compass mode &
\texttt{ap.mode="compass"} \\
Set wind mode &
\texttt{ap.mode="wind"} \\
Set true-wind mode &
\texttt{ap.mode="true wind"} \\
Set heading command &
\texttt{ap.heading\_command=123.4} \\
Select pilot &
\texttt{ap.pilot="basic"} \\
Watch values &
{\ttfamily\seqsplit{watch=\{"ap.heading":0.25,"ap.enabled":true\}}} \\
\bottomrule
\end{tabularx}
\end{table}
\subsection{Value Model}
The TCP protocol is organized around named values.  Each value has a name,
a current value, and metadata describing its type, writability,
persistence, range, units, or choices.  Values are arranged in dotted
namespaces.  For example,
\[
\texttt{ap.heading}
\]
is the current autopilot heading value, while
\[
\texttt{ap.heading\_command}
\]
is the commanded heading value.
The distinction between read-only values and writable properties is
central.  Sensor-like values, such as \texttt{ap.heading} or
\texttt{servo.current}, are normally read-only to remote TCP clients.
Control values, such as \texttt{ap.enabled}, \texttt{ap.mode}, and
\texttt{ap.heading\_command}, are writable properties.  When a remote
client sends a record for a writable value, the server validates the JSON
syntax and forwards the assignment to the owner of that value.  The owner
then applies value-specific validation such as range checking, enum-choice
checking, or boolean coercion.
The main value classes are summarized in
\cref{tab:pypilot-value-types}.
\begin{table}[!t]
\centering
\caption{Common pypilot value types exposed through the TCP protocol.}
\label{tab:pypilot-value-types}
\begin{tabularx}{\columnwidth}{Z{0.36\columnwidth}Y}
\toprule
Type & Meaning \\
\midrule
\texttt{Value} &
Generic value.  Not necessarily writable. \\
\texttt{SensorValue} &
Read-only or owner-updated measured value, often rounded for output. \\
\texttt{StringValue} &
String-valued state or diagnostic value. \\
{\ttfamily\seqsplit{BooleanProperty}} &
Writable boolean property. \\
\texttt{RangeProperty} &
Writable numeric property with minimum and maximum limits. \\
\texttt{RangeSetting} &
Writable range property with units and persistence. \\
\texttt{EnumProperty} &
Writable value constrained to a finite set of choices. \\
\texttt{ResettableValue} &
Writable resettable value, often reset by writing zero or false. \\
\bottomrule
\end{tabularx}
\end{table}
\subsection{Value Discovery}
A client can discover the available values by watching the special value
\texttt{values}.  The request is
\[
\texttt{watch=\{"values":true\}\textbackslash n}.
\]
The server replies with a line of the form
\[
\texttt{values=\{...\}\textbackslash n},
\]
where the right-hand side is a JSON object mapping value names to metadata
objects.  A representative fragment is
{\scriptsize
\[
\begin{aligned}
\texttt{values=\{}&
\texttt{"ap.enabled":\{"type":"BooleanProperty","writable":true\},}\\
&
\texttt{"ap.mode":\{"type":"EnumProperty","writable":true,}\\
&\quad
\texttt{"choices":["compass","gps","nav","wind","true wind"]\},}\\
&
\texttt{"ap.heading\_command":\{"type":"RangeProperty",}\\
&\quad
\texttt{"writable":true,"min":-180,"max":360\}\}}
\end{aligned}
\]
}
with the actual server response containing many more values.  This
metadata mechanism allows user interfaces to construct controls
automatically: check boxes for boolean properties, sliders for range
properties, selectors for enum properties, and read-only displays for
sensor values.
\subsection{Watch Mechanism}
The protocol is not purely request-response.  It is also a subscription
protocol.  A client subscribes to values by writing the special
\texttt{watch} value:
\[
\texttt{watch=\{"name":period\}\textbackslash n}.
\]
The watch period controls how the server sends updates.  A period of
\texttt{true} means continuous or change-driven watching.  A numeric
period means periodic updates with that period in seconds.  A value of
\texttt{false} removes the watch.
\begin{table}[!t]
\centering
\caption{Watch-period meanings.}
\label{tab:pypilot-watch-periods}
\begin{tabularx}{\columnwidth}{Z{0.36\columnwidth}Y}
\toprule
Period & Meaning \\
\midrule
\texttt{true} &
Watch continuously, equivalent internally to period \(0\). \\
\texttt{false} &
Unwatch the named value. \\
\texttt{0} &
Continuous watch. \\
positive number &
Periodic watch in seconds.  For example, \(0.25\) requests up to four
updates per second. \\
\bottomrule
\end{tabularx}
\end{table}
For example, a display that wants heading updates four times per second,
but only wants engagement and mode updates when they change, can send
\[
\begin{split}
\texttt{watch=\{}&
\texttt{"ap.heading":0.25,}\\
&
\texttt{"ap.heading\_command":true,}\\
&
\texttt{"ap.enabled":true,}\\
&
\texttt{"ap.mode":true\}\textbackslash n}.
\end{split}
\]
The server then emits records using the same assignment-line format:
\[
\begin{aligned}
&\texttt{ap.heading=277.3000}\\
&\texttt{ap.heading\_command=275.0}\\
&\texttt{ap.enabled=true}\\
&\texttt{ap.mode="compass"}.
\end{aligned}
\]
Therefore, after the watch is established, the same parser can be used for
both command responses and asynchronous updates.
\subsection{Autopilot Control Values}
The most important writable autopilot values are listed in
\cref{tab:pypilot-autopilot-values}.  These are the values an external
controller normally uses to engage the pilot, choose a steering mode, and
change the commanded heading.
\begin{table*}[!tbp]
\centering
\caption{Primary autopilot control values in the TCP protocol.}
\label{tab:pypilot-autopilot-values}
\begin{tabularx}{\textwidth}{Z{0.24\textwidth}Z{0.10\textwidth}Z{0.20\textwidth}Y}
\toprule
Value & Type & Example & Meaning \\
\midrule
\texttt{ap.enabled} &
{\ttfamily\seqsplit{BooleanProperty}} &
\texttt{ap.enabled=true} &
Engage or disengage the autopilot.  \texttt{false} is standby. \\
\texttt{ap.mode} &
\texttt{EnumProperty} &
\texttt{ap.mode="compass"} &
Select heading-reference mode.  Typical choices are
\texttt{"compass"}, \texttt{"gps"}, \texttt{"nav"},
\texttt{"wind"}, and \texttt{"true wind"}. \\
\texttt{ap.heading\_command} &
\texttt{RangeProperty} &
\texttt{ap.heading\_command=123.4} &
Set the commanded heading or wind-relative angle.  The value is resolved
according to the active mode and rounded internally to tenths of a degree. \\
\texttt{ap.pilot} &
\texttt{EnumProperty} &
\texttt{ap.pilot="basic"} &
Select the active pilot algorithm.  Available choices depend on the
loaded pilot modules. \\
\texttt{profile} &
\texttt{Value/Property} &
\texttt{profile="default"} &
Select the persistent pypilot profile.  Profiled gains and settings are
loaded from the selected profile. \\
\texttt{profiles} &
\texttt{Value/Property} &
\texttt{profiles=["default","heavy"]} &
Set or report the list of known profiles. \\
\bottomrule
\end{tabularx}
\end{table*}
A typical external controller sequence is:
\[
\begin{array}{ll}
1. & \text{connect to TCP port }23322,\\
2. & \text{watch } \texttt{values} \text{ or known status values},\\
3. & \text{set } \texttt{ap.mode},\\
4. & \text{set } \texttt{ap.heading\_command},\\
5. & \text{set } \texttt{ap.enabled=true},\\
6. & \text{adjust } \texttt{ap.heading\_command} \text{ as needed},\\
7. & \text{set } \texttt{ap.enabled=false} \text{ to stand by.}
\end{array}
\]
\subsection{Heading Command Semantics}
The meaning of \texttt{ap.heading\_command} depends on the active mode.
In compass, GPS, and navigation modes, it behaves as an absolute heading
command.  In wind modes, it behaves as a wind-relative command.  The
heading property class resolves values differently depending on whether
the active mode contains the word \texttt{wind}.  Conceptually,
\[
\begin{aligned}
\psi_c &=
\operatorname{resolv}(\psi_{\mathrm{input}},180^\circ)
&& \text{for compass/GPS/NAV},\\
\psi_c &=
\operatorname{resolv}(\psi_{\mathrm{input}},0^\circ)
&& \text{for wind modes}.
\end{aligned}
\]
The result is rounded to one decimal degree:
\[
\psi_c \leftarrow \frac{1}{10}
\operatorname{round}(10\psi_c).
\]
This means a client can send normal floating-point heading commands, but
the server-side autopilot state will use tenth-degree resolution.
\subsection{Errors and Invalid Requests}
Errors are also returned as assignment-style lines using the special name
\texttt{error}.  For example, if a client tries to write an unknown value,
the server can reply
\[
\texttt{error=invalid unknown value: some.name}.
\]
If a client tries to modify a value that exists but is not writable, the
server can reply
\[
\texttt{error=some.name is not writable}.
\]
If a line is syntactically malformed, lacks an equals sign, or contains an
invalid JSON right-hand side, the server reports an invalid request.  Empty
lines are ignored.  pypilot clients use this fact as a lightweight
keep-alive mechanism: after a quiet period, a client may send a blank
newline to make a broken TCP connection fail promptly.
\subsection{Optional UDP Output Stream}
The TCP protocol contains one special value for enabling UDP delivery of
watched output:
\[
\texttt{udp\_port=<port>}.
\]
After a TCP client sets \texttt{udp\_port} to a valid UDP port, the server
may deliver watched value updates to that UDP port on the same client IP
address.  The UDP datagram payload is still the same newline-delimited
text format:
\[
\begin{aligned}
&\texttt{ap.heading=123.4000}\\
&\texttt{ap.enabled=true}.
\end{aligned}
\]
The TCP connection remains the control channel.  UDP is an optimization
for watched output streams, especially small displays or remotes that want
fresh sensor values without accumulating TCP latency.  The watch
subscription itself is still configured with TCP records.
\subsection{Example Session}
The following example shows a complete minimal interaction.  The client
first subscribes to value metadata and a few status values, then commands
the pilot into compass mode at \(90^\circ\), and finally engages it.
\begin{quote}
\scriptsize\ttfamily
watch=\{"values":true\}\\
watch=\{"ap.enabled":true,"ap.mode":true,\allowbreak"ap.heading":0.5,\allowbreak"ap.heading\_command":true\}\\
ap.mode="compass"\\
ap.heading\_command=90.0\\
ap.enabled=true
\end{quote}
The server may respond with records such as
\begin{quote}
\scriptsize\ttfamily
values=\{"ap.enabled":\{"type":"BooleanProperty",\allowbreak"writable":true\},...\}\\
ap.mode="compass"\\
ap.heading\_command=90.0\\
ap.enabled=true\\
ap.heading=88.7000\\
ap.heading=88.9000
\end{quote}
To stand by, the client sends
\begin{quote}
\ttfamily
ap.enabled=false
\end{quote}
\subsection{Comparison with the NMEA TCP Server}
pypilot also contains a separate NMEA TCP interface, normally on port
\[
p_{\mathrm{nmea}} = 20220.
\]
That port carries ordinary NMEA sentences such as GPS, wind, rudder, and
APB/navigation data:
\[
\texttt{\$...*hh\textbackslash r\textbackslash n}.
\]
The NMEA TCP server is therefore a marine-data interface.  It can
influence the autopilot indirectly by supplying GPS, wind, or APB
cross-track-error data, especially in GPS and NAV modes.  It is not the
native pypilot command/value protocol.
The distinction is summarized in
\cref{tab:pypilot-tcp-vs-nmea}.
\begin{table}[!t]
\centering
\caption{pypilot command TCP versus NMEA TCP.}
\label{tab:pypilot-tcp-vs-nmea}
\begin{tabularx}{\columnwidth}{Z{0.20\columnwidth}YY}
\toprule
Field & pypilot TCP & NMEA TCP \\
\midrule
Default port & \texttt{23322} & \texttt{20220} \\
Record type & \texttt{name=json} & NMEA sentence \\
Primary use & Control/status/config & Marine sensor/navigation data \\
Command example & \texttt{ap.enabled=true} & APB/RMC/MWV sentence \\
Framing & \texttt{\textbackslash n} & \texttt{\textbackslash r\textbackslash n} \\
Self-description & \texttt{values} metadata & none \\
Subscriptions & \texttt{watch=\{...\}} & not native \\
\bottomrule
\end{tabularx}
\end{table}
For direct autopilot control, clients should use the pypilot TCP value
protocol on port \texttt{23322}.  For navigation and sensor sentences,
clients should use the NMEA TCP server on port \texttt{20220}.
\subsection{Implementation Notes for External Clients}
An external client implementation only needs a line-oriented TCP parser.
The minimum algorithm is:
\begin{enumerate}
\setlength{\itemsep}{0pt}
\setlength{\parskip}{0pt}
\item Open a TCP connection to \texttt{host:23322}.
\item Send newline-terminated assignment records.
\item Split incoming lines at the first \texttt{=}.
\item Decode the right side as JSON.
\item Handle \texttt{error} records separately.
\item Use \texttt{watch} for asynchronous updates.
\end{enumerate}

A minimal parser can be written around the following rule:
\[
(\text{name},\text{payload})
=
\operatorname{split}_{1}(\text{line},\texttt{=}).
\]
Then
\[
\text{value} = \operatorname{JSONDecode}(\text{payload}).
\]
If \(\text{name}=\texttt{error}\), the payload should be treated as a
diagnostic string rather than a normal value update.
Because updates are asynchronous, a client should not assume that the next
line received is a direct response to the last command sent.  After watches
are installed, any watched value can be emitted whenever its watch policy
allows it.  The protocol is therefore best understood as a shared
publish/subscribe value bus with writable properties, rather than as a
strict request-response API.

\section{Control Parameters, Gates, and Filters}
\label{sec:pypilot-parameters}
The pypilot control stack contains many explicit gains, filters, clamps,
timeouts, and gating rules.  These are not incidental constants; they are
part of the practical behavior of the autopilot.  The main active
subsystems are the pilot controller, the heading and inertial filters, the
wind and navigation filters, the tacking state machine, and the servo
driver.  This section summarizes the important parameters found in the
source code, their defaults, their code-enforced ranges where applicable,
and reasonable practical tuning ranges.
Throughout the code, many filters are first-order exponential moving
averages of the form
\[
x_f[k]
=
(1-\alpha)x_f[k-1] + \alpha x[k],
\]
where \(\alpha\) is the filter constant.  For small \(\alpha\), the
approximate time constant is
\[
\tau \approx \frac{\Delta t}{\alpha}.
\]
Thus, at the default IMU/control rate of \SI{20}{Hz}, where
\(\Delta t = \SI{0.05}{s}\), a filter constant of \(\alpha=0.2\)
corresponds to a fast time constant of roughly \(\SI{0.25}{s}\), while
\(\alpha=0.01\) corresponds to a much slower filter.
The tables below distinguish between two different notions of range:
\emph{code range} is the range accepted by the pypilot value system, while
\emph{practical range} is a conservative operating range for real tuning.
The practical ranges are intentionally narrower than the code ranges.
\subsection{Active Pilot Gains}
In the inspected source tree, the normally active pilot modules are
\texttt{basic}, \texttt{wind}, and \texttt{absolute}.  Several other
pilots exist in the source tree, but are marked disabled and should be
treated as experimental unless explicitly enabled.
\subsubsection{Basic Pilot}
The basic pilot is the normal compass/GPS/NAV heading-hold controller.  It
uses an extended proportional--derivative form:
\[
u
=
K_P e
+
K_D \dot{\psi}
+
K_{DD} \ddot{\psi}
+
K_{PR}\operatorname{sgn}(e)\sqrt{\abs{e}}
+
K_{FF}\dot{\psi}_c ,
\]
where \(u\) is the normalized servo command, \(e\) is the heading error,
\(\dot{\psi}\) is measured heading rate, \(\ddot{\psi}\) is heading
rate-of-rate, and \(\dot{\psi}_c\) is the filtered change in commanded
heading.  The square-root term is unusual compared with a conventional PID
controller.  It gives useful authority near small and moderate heading
errors without requiring a large linear proportional gain.
\begin{table*}[!tbp]
\centering
\caption{Active pilot gains for the basic, wind, and absolute pilots.}
\label{tab:pypilot-active-pilot-gains}
\label{tab:pypilot-basic-gains}
\label{tab:pypilot-wind-gains}
\label{tab:pypilot-absolute-gains}
\begin{tabularx}{\textwidth}{Z{0.13\textwidth}Z{0.20\textwidth}c c Z{0.13\textwidth}Y}
\toprule
Pilot & Value & Default & Code range & Practical range & Meaning \\
\midrule
\multirow{5}{*}{\texttt{basic}} & \texttt{ap.pilot.basic.P} &
\num{0.003} & \(\num{0}\ldots\num{0.03}\) & \(\num{0.001}\ldots\num{0.010}\) &
Linear heading-error gain.  Higher values turn harder toward the commanded heading; excessive values cause hunting. \\
& \texttt{ap.pilot.basic.D} &
\num{0.09} & \(\num{0}\ldots\num{0.24}\) & \(\num{0.04}\ldots\num{0.16}\) &
Heading-rate damping from the gyro.  Higher values suppress oscillation but can make the response sluggish. \\
& \texttt{ap.pilot.basic.DD} &
\num{0.075} & \(\num{0}\ldots\num{0.24}\) & \(\num{0}\ldots\num{0.14}\) &
Heading-acceleration damping.  Useful for suppressing sharp yaw-rate changes; too much can amplify noise. \\
& \texttt{ap.pilot.basic.PR} &
\num{0.005} & \(\num{0}\ldots\num{0.02}\) & \(\num{0.002}\ldots\num{0.010}\) &
Square-root heading-error gain.  Adds authority without requiring a large linear \(P\) gain. \\
& \texttt{ap.pilot.basic.FF} &
\num{0.6} & \(\num{0}\ldots\num{2.4}\) & \(\num{0.2}\ldots\num{1.2}\) &
Feed-forward from commanded heading changes.  Improves response to button pushes or external course changes. \\
\midrule
\multirow{4}{*}{\texttt{wind}} & \texttt{ap.pilot.wind.P} &
\num{0.003} & \(\num{0}\ldots\num{0.02}\) & \(\num{0.001}\ldots\num{0.008}\) &
Wind-angle or wind-referenced heading-error gain. \\
& \texttt{ap.pilot.wind.D} &
\num{0.1} & \(\num{0}\ldots\num{1.0}\) & \(\num{0.04}\ldots\num{0.25}\) &
Heading-rate damping. \\
& \texttt{ap.pilot.wind.DD} &
\num{0.05} & \(\num{0}\ldots\num{1.0}\) & \(\num{0}\ldots\num{0.15}\) &
Heading-acceleration damping. \\
& \texttt{ap.pilot.wind.WG} &
\num{0} & \(\num{-0.1}\ldots\num{0.1}\) & \(\num{-0.03}\ldots\num{0.03}\) &
Wind-gust compensation.  Should normally remain near zero unless wind data are clean and well-filtered. \\
\midrule
\multirow{4}{*}{\texttt{absolute}} & \texttt{ap.pilot.absolute.P} &
\num{0.05} & \(\num{0}\ldots\num{0.3}\) & \(\num{0.02}\ldots\num{0.12}\) &
Converts heading error into rudder-position command. \\
& \texttt{ap.pilot.absolute.I} &
\num{0} & \(\num{0}\ldots\num{0.8}\) & \(\num{0}\ldots\num{0.05}\) &
Integral correction.  Can wind up if used aggressively. \\
& \texttt{ap.pilot.absolute.D} &
\num{0.05} & \(\num{0}\ldots\num{0.5}\) & \(\num{0.02}\ldots\num{0.15}\) &
Heading-rate damping. \\
& \texttt{ap.pilot.absolute.FF} &
\num{0} & \(\num{0}\ldots\num{1.0}\) & \(\num{0}\ldots\num{0.5}\) &
Feed-forward from heading-command changes. \\
\bottomrule
\end{tabularx}
\end{table*}
Practical symptoms and adjustments are summarized in
\cref{tab:pypilot-basic-tuning}.
\begin{table}[!t]
\centering
\caption{Basic pilot tuning symptoms.}
\label{tab:pypilot-basic-tuning}
\begin{tabularx}{\columnwidth}{Z{0.36\columnwidth}Y}
\toprule
Observed behavior & Likely adjustment \\
\midrule
Slowly wanders off course &
Increase \texttt{P} or \texttt{PR}. \\
Oscillates around course &
Decrease \texttt{P}/\texttt{PR}, or increase \texttt{D}. \\
Rudder chatters or reacts to noisy yaw acceleration &
Decrease \texttt{DD}. \\
Course changes feel delayed &
Increase \texttt{FF}. \\
Course changes overshoot badly &
Decrease \texttt{FF} or \texttt{P}; increase \texttt{D}. \\
\bottomrule
\end{tabularx}
\end{table}
\subsubsection{Wind Pilot}
The wind pilot has a similar structure but adds a wind-gust term:
\[
u
=
K_P e
+
K_D \dot{\psi}
+
K_{DD} \ddot{\psi}
+
K_{WG} g_w ,
\]
where \(g_w\) is the apparent wind-speed change used as a gust signal.
The \texttt{WG} gain should be treated carefully.  A noisy wind sensor will
inject noise directly into the rudder command if \texttt{WG} is increased.
\subsubsection{Absolute Pilot}
The absolute pilot attempts to command a rudder position rather than a
relative motor speed.  It requires valid rudder feedback.  In the inspected
source, it is present and active, but should be regarded as experimental
because its behavior depends strongly on rudder calibration.
\subsection{Heading Error, Command Rate, and Integral Gates}
The autopilot does not feed unbounded heading error into the pilot.  The
heading error is explicitly clamped:
\[
e
=
\operatorname{clamp}
\left(
\operatorname{resolv}(\psi-\psi_c),
-30^\circ,
+30^\circ
\right).
\]
Therefore, even if the vessel is far off course, the control law sees at
most \(30^\circ\) of error.  This is a practical marine safety measure:
it prevents extreme actuator commands after a sensor jump, a large
disturbance, a tack, or a mode transition.
\begin{table}[!t]
\centering
\caption{Heading-error and command-rate gates.}
\label{tab:pypilot-heading-gates}
\begin{tabularx}{\columnwidth}{Z{0.26\columnwidth}cZ{0.14\columnwidth}Y}
\toprule
Gate/filter & Value & Tunable & Meaning \\
\midrule
Heading-error clamp &
\(\pm 30^\circ\) &
No &
Maximum error seen by the pilot controller. \\
Heading-command resolution &
\(\num{0.1}^\circ\) &
No &
Commanded heading is rounded to tenths of a degree. \\
Integral \(dt\) cap &
\(\SI{1}{s}\) &
No &
Prevents a scheduler stall from producing a large integral jump. \\
Integral clamp &
\(\pm 10\) &
No &
Limits stored heading-error integral. \\
Feed-forward EMA alpha &
\num{0.1} &
No &
Smooths heading-command changes before use by \texttt{FF}. \\
Feed-forward reset timeout &
\(\SI{1}{s}\) &
No &
If the last command is older than this, the feed-forward reference resets. \\
\bottomrule
\end{tabularx}
\end{table}
The heading-error integral is computed as
\[
I_e[k]
=
\operatorname{clamp}
\left(
I_e[k-1]
+
\frac{e[k]}{50}\Delta t,
-10,
10
\right).
\]
The active \texttt{basic} and \texttt{wind} pilots do not use this
integral term in their default equations.  It is more relevant to the
\texttt{absolute} pilot and to experimental or future pilots.
The feed-forward command-rate signal is updated with
\[
\dot{\psi}_{c,f}[k]
=
0.9\,\dot{\psi}_{c,f}[k-1]
+
0.1\,\Delta\psi_c[k].
\]
In non-wind modes, the sign is inverted before filtering; wind modes use
the opposite sign convention because apparent wind direction is a
``direction from'' measurement.
\subsection{IMU and Heading Filters}
The default IMU/control loop rate is \(\SI{20}{Hz}\).  The allowed rate
choices are \(\SI{10}{Hz}\) and \(\SI{20}{Hz}\).
\begin{table*}[!tbp]
\centering
\caption{IMU and heading filter parameters.}
\label{tab:pypilot-imu-filters}
\begin{tabularx}{\textwidth}{Z{0.23\textwidth}c c Z{0.16\textwidth}Y}
\toprule
Value & Default & Code range & Practical range & Meaning \\
\midrule
\texttt{imu.rate} &
\(\SI{20}{Hz}\) &
\(\SI{10}{Hz}\) or \(\SI{20}{Hz}\) &
\(\SI{20}{Hz}\) &
Main IMU and autopilot loop rate. \\
\texttt{imu.heading\_offset} &
\(90^\circ\) &
\(-180^\circ\ldots180^\circ\) &
installation-dependent &
Compass/IMU heading offset.  Default reflects the expected display-facing mounting orientation. \\
\texttt{imu.heading\_lowpass\_constant} &
\num{0.2} &
\(\num{0.05}\ldots\num{0.3}\) &
\(\num{0.1}\ldots\num{0.25}\) &
EMA alpha for heading. \\
\texttt{imu.headingrate\_lowpass\_constant} &
\num{0.2} &
\(\num{0.05}\ldots\num{0.3}\) &
\(\num{0.1}\ldots\num{0.25}\) &
EMA alpha for gyro-derived heading rate. \\
{\ttfamily\seqsplit{imu.headingraterate\_lowpass\_constant}} &
\num{0.1} &
\(\num{0.05}\ldots\num{0.3}\) &
\(\num{0.05}\ldots\num{0.15}\) &
EMA alpha for heading acceleration. \\
\bottomrule
\end{tabularx}
\end{table*}
At \(\SI{20}{Hz}\), the approximate time constants are shown in
\cref{tab:pypilot-alpha-timeconstants}.
\begin{table}[!t]
\centering
\caption{Approximate EMA time constants at \(\SI{20}{Hz}\).}
\label{tab:pypilot-alpha-timeconstants}
\begin{tabular}{cc}
\toprule
\(\alpha\) & Approximate time constant \\
\midrule
\num{0.05} & \(\SI{1.0}{s}\) \\
\num{0.10} & \(\SI{0.5}{s}\) \\
\num{0.20} & \(\SI{0.25}{s}\) \\
\num{0.30} & \(\SI{0.17}{s}\) \\
\bottomrule
\end{tabular}
\end{table}
The heading rate-of-rate estimate is guarded by an IMU timing gate:
\[
\ddot{\psi}
=
\begin{cases}
\dfrac{\dot{\psi}[k]-\dot{\psi}[k-1]}{\Delta t},
& \num{0.01} < \Delta t < \num{0.2},\\[8pt]
0, & \text{otherwise}.
\end{cases}
\]
This prevents malformed timing intervals from producing large
\(\ddot{\psi}\) spikes.
The heel estimate is another fixed EMA:
\[
\phi_{\mathrm{heel}}[k]
=
0.97\,\phi_{\mathrm{heel}}[k-1]
+
0.03\,\phi_{\mathrm{roll}}[k].
\]
At \(\SI{20}{Hz}\), this is a slow filter with a time constant of roughly
\(\SI{1.5}{s}\) to \(\SI{1.7}{s}\).
The IMU warning logic also includes a simple alignment gate:
\[
\abs{\theta_{\mathrm{pitch}}} > 35^\circ
\quad\text{or}\quad
\abs{\theta_{\mathrm{roll}}} > 35^\circ .
\]
If either condition is true, pypilot appends an alignment warning.
\subsection{Compass, GPS, and Wind Offset Filters}
pypilot maintains long-term offsets between compass heading, GPS track,
apparent wind, and true wind.  These offsets are not direct steering
commands; they are slowly adapted references used to preserve continuity
between modes.
The GPS speed used for offset weighting is filtered by
\[
V_{g,f}[k]
=
0.998\,V_{g,f}[k-1]
+
0.002\,V_g[k].
\]
The GPS-to-compass offset is only updated above a speed gate:
\[
V_g > \SI{1}{kn}.
\]
When this gate is satisfied, the update factor is
\[
\alpha_g
=
0.005\log(V_{g,f}+1).
\]
This gives GPS track more authority at higher speed and avoids using GPS
course-over-ground when the vessel is nearly stationary.
\begin{table*}[!tbp]
\centering
\caption{Compass, GPS, and wind offset filters.}
\label{tab:pypilot-offset-filters}
\begin{tabularx}{\textwidth}{Z{0.23\textwidth}c c Z{0.16\textwidth}Y}
\toprule
Quantity & Default/fixed value & Code range & Practical range & Meaning \\
\midrule
GPS speed EMA alpha &
\num{0.002} &
fixed &
fixed &
Very slow smoothing of GPS speed before GPS/compass offset updates. \\
GPS offset speed gate &
\(\SI{1}{kn}\) &
fixed &
\(\SI{1}{kn}\ldots\SI{3}{kn}\) if editing source &
Avoids trusting GPS track at very low speed. \\
GPS offset update alpha &
\(0.005\log(V_{g,f}+1)\) &
fixed &
fixed &
Weights GPS/compass offset update more heavily at higher speed. \\
\texttt{ap.wind\_offset\_filter} &
\num{0.1} &
\(\num{0.01}\ldots\num{0.5}\) &
\(\num{0.03}\ldots\num{0.2}\) &
EMA alpha for apparent-wind and true-wind compass offsets. \\
\bottomrule
\end{tabularx}
\end{table*}
The wind offset update has the form
\[
\beta_f[k]
=
(1-\alpha_w)\beta_f[k-1]
+
\alpha_w \beta[k],
\]
where the user-facing parameter is
\[
\alpha_w = \texttt{ap.wind\_offset\_filter}.
\]
Large values adapt quickly to wind shifts but can chase wind sensor noise.
Small values preserve a stable long-term wind/compass offset.
\subsection{Wind Sensor Filters}
The apparent-wind and true-wind sensor paths use both fixed and tunable
filters.  The base wind sensor rate defaults to \(\SI{4}{Hz}\), and the
wind speed filter is deliberately slow:
\[
V_{w,f}[k]
=
0.99\,V_{w,f}[k-1]
+
0.01\,V_w[k].
\]
At \(\SI{4}{Hz}\), the approximate time constant is about
\(\SI{25}{s}\).  This slow speed estimate is then used to weight the
direction filter.
\begin{table*}[!tbp]
\centering
\caption{Wind and true-wind filtering parameters.}
\label{tab:pypilot-wind-filters}
\begin{tabularx}{\textwidth}{Z{0.23\textwidth}c c Z{0.16\textwidth}Y}
\toprule
Value & Default & Code range & Practical range & Meaning \\
\midrule
\texttt{wind.rate} &
\(\SI{4}{Hz}\) &
\(\SI{0}{Hz}\ldots\SI{50}{Hz}\) &
\(\SI{2}{Hz}\ldots\SI{10}{Hz}\) &
Requested/input apparent-wind data rate. \\
\texttt{truewind.rate} &
\(\SI{4}{Hz}\) &
\(\SI{0}{Hz}\ldots\SI{50}{Hz}\) &
\(\SI{2}{Hz}\ldots\SI{10}{Hz}\) &
Requested/input true-wind data rate. \\
\texttt{wind.offset} &
\(0^\circ\) &
\(-180^\circ\ldots180^\circ\) &
calibrated value &
Apparent-wind angle offset. \\
\texttt{truewind.offset} &
\(0^\circ\) &
\(-180^\circ\ldots180^\circ\) &
calibrated value &
True-wind angle offset if true wind is externally supplied. \\
\texttt{wind.sensors\_height} &
\(\SI{0}{m}\) &
\(\SI{0}{m}\ldots\SI{100}{m}\) &
actual sensor height &
Height used for mast-motion apparent-wind compensation. \\
\texttt{wind.filter\_constant} &
\num{0.1} &
\(\num{0.01}\ldots\num{1.0}\) &
\(\num{0.03}\ldots\num{0.25}\) &
Base coefficient for speed-weighted wind-direction filtering. \\
\texttt{truewind.filter\_constant} &
\num{0.1} &
\(\num{0.01}\ldots\num{1.0}\) &
\(\num{0.03}\ldots\num{0.25}\) &
Same as above for the true-wind path. \\
\bottomrule
\end{tabularx}
\end{table*}
The wind-direction filter coefficient is not constant.  It is computed as
\[
\alpha_{\mathrm{wind}}
=
c_w
\log\left(\frac{V_{w,f}}{5}+1.1\right),
\]
where
\[
c_w = \texttt{wind.filter\_constant}.
\]
The filtered wind direction is then updated by
\[
\theta_{w,f}[k]
=
(1-\alpha_{\mathrm{wind}})\theta_{w,f}[k-1]
+
\alpha_{\mathrm{wind}}\theta_w[k],
\]
with directional wrap handled by the angle-resolution function.
With the default \(c_w=0.1\), the effective alpha is approximately
\cref{tab:pypilot-wind-alpha}.
\begin{table}[!t]
\centering
\caption{Approximate wind-direction alpha for default \texttt{filter\_constant}.}
\label{tab:pypilot-wind-alpha}
\begin{tabular}{cc}
\toprule
Filtered wind speed & Effective \(\alpha\) \\
\midrule
\(\SI{0}{kn}\)  & \(\num{0.0095}\) \\
\(\SI{5}{kn}\)  & \(\num{0.074}\) \\
\(\SI{10}{kn}\) & \(\num{0.113}\) \\
\(\SI{20}{kn}\) & \(\num{0.163}\) \\
\bottomrule
\end{tabular}
\end{table}
Since the code does not clamp this computed alpha, large
\texttt{filter\_constant} values can produce \(\alpha>1\) at high wind
speed.  For robust behavior, \texttt{filter\_constant} should normally be
kept well below \(1\).
If \texttt{wind.sensors\_height} is nonzero, apparent-wind measurements
are compensated using pitch and roll rates.  The source computes a mast
motion correction approximately proportional to
\[
\kappa
=
\frac{\pi}{180}
\cdot
1.94
\cdot
h,
\]
where \(h\) is the sensor height in meters and the factor \(1.94\)
converts meters per second to knots.  The apparent wind vector is adjusted
by roll-rate and pitch-rate terms.  This should only be enabled when IMU
alignment and sensor geometry are trusted.
\subsection{GPS, APB, and Navigation Parameters}
The GPS sensor defaults to \(\SI{1}{Hz}\), while the APB navigation
interface contains the main cross-track-error gain used in NAV mode.
\begin{table*}[!tbp]
\centering
\caption{GPS, APB navigation, and user-visible GPS filter values.}
\label{tab:pypilot-gps-apb-filter-parameters}
\begin{tabularx}{\textwidth}{Z{0.23\textwidth}c c Z{0.16\textwidth}Y}
\toprule
Value/gate & Default & Code range & Practical range & Meaning \\
\midrule
\multicolumn{5}{l}{\textit{GPS and APB navigation parameters}} \\
\addlinespace[0.25ex]
\texttt{gps.rate} &
\(\SI{1}{Hz}\) &
\(\SI{0}{Hz}\ldots\SI{50}{Hz}\) &
\(\SI{1}{Hz}\ldots\SI{10}{Hz}\) &
Requested/input GPS update rate. \\
\texttt{apb.xte.gain} &
\num{300} &
\(\num{0}\ldots\num{3000}\) &
\(\num{50}\ldots\num{600}\) &
Cross-track-error gain used in NAV mode. \\
APB accept interval &
\(\SI{0.5}{s}\) &
fixed &
fixed &
APB updates faster than \(\SI{2}{Hz}\) are ignored. \\
Heading-command update threshold &
\(0.1^\circ\) &
fixed &
fixed &
NAV mode only writes a new command if it differs by more than this. \\
\addlinespace[0.6ex]
\multicolumn{5}{l}{\textit{User-visible GPS filter values}} \\
\addlinespace[0.25ex]
\texttt{gps.filtered.enabled} &
\texttt{false} &
boolean &
usually \texttt{false} &
Enables the GPS/IMU filter process. \\
\texttt{gps.filtered.output} &
\texttt{false} &
boolean &
usually \texttt{false} &
Controls whether filtered GPS output is used/exported. \\
\texttt{gps.filtered.time\_offset} &
\(\SI{0.7}{s}\) &
not range-limited &
\(\SI{0.3}{s}\ldots\SI{1.5}{s}\) &
Assumed GPS latency. \\
\addlinespace[0.6ex]
\multicolumn{5}{l}{\textit{Fixed GPS filter constants}} \\
\addlinespace[0.25ex]
Position measurement sigma & \(\SI{10}{m}\) & fixed & fixed & GPS position noise assumption. \\
Velocity measurement sigma & \(\SI{0.25}{m/s}\) & fixed & fixed & GPS velocity noise assumption. \\
Initial position deviation & \(\SI{30}{m}\) & fixed & fixed & Initial/process covariance scale. \\
Initial velocity deviation & \(\SI{3}{m/s}\) & fixed & fixed & Initial/process covariance scale. \\
3D mode & \texttt{false} & fixed & fixed & The filter runs in 2D by default. \\
\bottomrule
\end{tabularx}
\end{table*}
The NAV-mode command is very simple:
\[
\psi_c
=
\psi_{\mathrm{track}}
+
K_{\mathrm{xte}}\,e_{\mathrm{xte}},
\]
where
\[
K_{\mathrm{xte}} = \texttt{apb.xte.gain}.
\]
The source comment states that the default value \(K_{\mathrm{xte}}=300\)
corresponds to a \(30^\circ\) correction for \(0.1\) mile of cross-track
error.  This is a proportional cross-track correction, not an L1 or
pure-pursuit path-following law.
The APB update only affects the heading command when
\[
\texttt{ap.mode}=\texttt{"nav"},
\qquad
\texttt{ap.enabled}=\texttt{true},
\]
and the APB data mode is accepted as GPS-derived navigation data.
The optional filtered GPS process is disabled by default.  It estimates
position and velocity with a simple Kalman structure using GPS fixes and
IMU acceleration.  It is not the normal steering loop in the inspected
configuration.
The GPS filter also contains several timing gates:
\[
\Delta t_{\mathrm{predict}} < 0
\quad\text{or}\quad
\Delta t_{\mathrm{predict}} > \SI{0.5}{s}
\quad\Longrightarrow\quad
\text{reset}.
\]
For normal predictions, the applied step is clamped:
\[
\Delta t
\leftarrow
\operatorname{clamp}(\Delta t,\SI{0.02}{s},\SI{0.1}{s}).
\]
GPS fixes older than \(\SI{5}{s}\) are considered stale; after more than
five stale fixes, the system time offset and filter are reset.  Magnetic
declination is refreshed from the magnetic model no more often than once
every four hours.
\subsection{Tacking Parameters and Gates}
The tacking subsystem is a small state machine with states
\[
\texttt{none},\quad
\texttt{begin},\quad
\texttt{waiting},\quad
\texttt{tacking}.
\]
External clients initiate a tack by setting the state to
\texttt{begin}.  The state machine then waits for the configured delay,
commands a turn, and finishes the tack when the configured threshold is
crossed.
\begin{table*}[!tbp]
\centering
\caption{Tacking parameters.}
\label{tab:pypilot-tacking-parameters}
\begin{tabularx}{\textwidth}{Z{0.23\textwidth}c c Z{0.16\textwidth}Y}
\toprule
Value & Default & Code range & Practical range & Meaning \\
\midrule
\multicolumn{5}{l}{\textit{Configurable tacking parameters}} \\
\addlinespace[0.25ex]
\texttt{ap.tack.delay} &
\(\SI{0}{s}\) &
\(\SI{0}{s}\ldots\SI{60}{s}\) &
\(\SI{0}{s}\ldots\SI{10}{s}\) &
Delay between tack request and start of tack. \\
\texttt{ap.tack.angle} &
\(100^\circ\) &
\(10^\circ\ldots180^\circ\) &
\(70^\circ\ldots120^\circ\) &
Nominal tack angle in compass mode. \\
\texttt{ap.tack.rate} &
\(\SI{15}{deg/s}\) &
\(\SI{1}{deg/s}\ldots\SI{100}{deg/s}\) &
\(\SI{5}{deg/s}\ldots\SI{25}{deg/s}\) &
Turn-rate scaling during the tack. \\
\texttt{ap.tack.threshold} &
\(\SI{50}{\percent}\) &
\(\SI{10}{\percent}\ldots\SI{100}{\percent}\) &
\(\SI{40}{\percent}\ldots\SI{70}{\percent}\) &
Percent completion at which the tack is considered complete. \\
\texttt{ap.tack.use\_heel} &
\texttt{false} &
boolean &
usually \texttt{false} &
Infer tack side from heel. \\
\texttt{ap.tack.use\_wind\_direction} &
\texttt{true} &
boolean &
usually \texttt{true} &
Infer tack side from apparent wind direction. \\
\addlinespace[0.6ex]
\multicolumn{5}{l}{\textit{Fixed tacking detector gates}} \\
\addlinespace[0.25ex]
Direction-log update interval & \(\SI{0.25}{s}\) & fixed & fixed & Tack-direction logs are updated at most this often. \\
Direction-log clear interval & \(\SI{1}{s}\) & fixed & fixed & Longer update gaps clear the tack-direction log. \\
Apparent-wind detector threshold & \(20^\circ\) & fixed & fixed & Apparent-wind side-change detector threshold. \\
Heel detector threshold & \(7^\circ\) & fixed & fixed & Heel side-change detector threshold. \\
Heel detector warmup & \(\SI{30}{s}\) & fixed & fixed & Heel is ignored until this delay has elapsed. \\
\bottomrule
\end{tabularx}
\end{table*}
The tack direction detector maintains logs that are updated at most every
\(\SI{0.25}{s}\).  If the update interval exceeds \(\SI{1}{s}\), the log
is cleared.  The apparent-wind tack detector uses a threshold of
\(20^\circ\), and the heel detector uses a threshold of \(7^\circ\).  Heel
is not used immediately; the code waits more than \(\SI{30}{s}\) before
using heel as a tack-side detector.
During the tacking state, the command is computed from turn-rate terms:
\[
u_{\mathrm{tack}}
=
\frac{\dot{\psi}+\ddot{\psi}/2}
{\texttt{ap.tack.rate}}
+
d,
\]
where \(d\in\{-1,+1\}\) is the tack direction.  The result is clamped:
\[
u_{\mathrm{tack}}
\leftarrow
\operatorname{clamp}(u_{\mathrm{tack}},-1,+1).
\]
This prevents excessive command values during the tack.
\subsection{Servo and Actuator Parameters}
The servo layer converts the pilot's normalized command into motor drive,
protects the actuator, handles current and temperature limits, and shapes
small commands into practical pulses.  This layer is especially important
because pypilot is designed to work with real DC actuators and motor
controllers rather than idealized rudder servos.
\subsubsection{Safety Limits}
\begin{table*}[!tbp]
\centering
\caption{Servo safety limits.}
\label{tab:pypilot-servo-safety}
\begin{tabularx}{\textwidth}{Z{0.23\textwidth}c c Z{0.16\textwidth}Y}
\toprule
Value & Default & Code range & Practical range & Meaning \\
\midrule
\texttt{servo.max\_current} &
\(\SI{4.5}{A}\) &
\(\SI{0}{A}\ldots\SI{50}{A}\) &
drive-dependent &
Motor current trip limit sent to the controller. \\
\texttt{servo.max\_controller\_temp} &
\(\SI{60}{\celsius}\) &
\(\SI{45}{\celsius}\ldots\SI{80}{\celsius}\) &
\(\SI{55}{\celsius}\ldots\SI{70}{\celsius}\) &
Controller over-temperature limit. \\
\texttt{servo.max\_motor\_temp} &
\(\SI{60}{\celsius}\) &
\(\SI{30}{\celsius}\ldots\SI{80}{\celsius}\) &
\(\SI{50}{\celsius}\ldots\SI{70}{\celsius}\) &
Motor over-temperature limit. \\
\texttt{servo.faults} &
\num{0} &
resettable &
diagnostic &
Persistent fault counter. \\
\bottomrule
\end{tabularx}
\end{table*}
The practical current limit depends on the actuator.  Small tiller pilots
may use only a few amperes, while larger hydraulic or wheel drives can
require much higher limits.  The Arduino controller also reports its
current range, and the pypilot host adjusts the allowed maximum
accordingly.
\subsubsection{Servo Output Scaling}
\begin{table*}[!tbp]
\centering
\caption{Servo output scaling and pulse-shaping parameters.}
\label{tab:pypilot-servo-scaling}
\begin{tabularx}{\textwidth}{Z{0.23\textwidth}c c Z{0.16\textwidth}Y}
\toprule
Value & Default & Code range & Practical range & Meaning \\
\midrule
\texttt{servo.gain} &
\num{1} &
\(\num{-10}\ldots\num{10}\) &
\(\num{0.5}\ldots\num{2}\), or \(-1\) if reversed &
Multiplies final speed command.  Negative values reverse the drive direction. \\
\texttt{servo.speed.min} &
\(\SI{100}{\percent}\) &
\(\SI{0}{\percent}\ldots\SI{100}{\percent}\) &
\(\SI{20}{\percent}\ldots\SI{60}{\percent}\) after calibration &
Minimum motor speed once a pulse is emitted. \\
\texttt{servo.speed.max} &
\(\SI{100}{\percent}\) &
\(\SI{0}{\percent}\ldots\SI{100}{\percent}\) &
\(\SI{70}{\percent}\ldots\SI{100}{\percent}\) &
Maximum motor speed. \\
\texttt{servo.speed\_gain} &
\num{0} &
\(\num{0}\ldots\num{1}\) &
\(\num{0}\ldots\num{0.3}\) &
Raises minimum speed as duty cycle increases. \\
\texttt{servo.period} &
\(\SI{0.4}{s}\) &
\(\SI{0.1}{s}\ldots\SI{1.0}{s}\) &
\(\SI{0.25}{s}\ldots\SI{0.6}{s}\) &
Pulse/windup period used to avoid tiny actuator movements. \\
\bottomrule
\end{tabularx}
\end{table*}
The source enforces
\[
\texttt{servo.speed.max}
\ge
\texttt{servo.speed.min}.
\]
The default \(\SI{100}{\percent}\) minimum and maximum produce a
full-speed pulsed actuator.  After calibration, lowering the minimum speed
can give smoother steering, but it must remain high enough to reliably
move the actuator.
\subsubsection{Servo Windup and Pulse Logic}
Normal autopilot commands are not always sent directly to the motor.  The
servo layer integrates a windup term:
\[
w[k]
=
w[k-1]
+
\left(u[k]-u_{\mathrm{actual}}[k]\right)\Delta t .
\]
A motor pulse is emitted only when
\[
\abs{w}
>
\frac{T_s u_{\min}}{1.5},
\]
where
\[
T_s = \texttt{servo.period}
\]
and \(u_{\min}\) is the configured minimum speed.  The windup is bounded:
\[
\abs{w}
\le
1.5T_s .
\]
If the windup saturates, pypilot sets an internal saturation flag.
The servo also suppresses rapid direction changes.  If the requested
direction changes within one \texttt{servo.period}, the previous direction
or a zero command may be retained briefly.  This protects the actuator
from rapid reversals and prevents tiny alternating commands from wearing
the drive.
\subsubsection{Manual Override and Command Timeout}
Manual commands and autopilot-generated commands are treated differently.
External/manual writes to \texttt{servo.command} are immediate, while
autopilot-generated calls use the period/windup logic.
\begin{table}[!t]
\centering
\caption{Servo command timing gates.}
\label{tab:pypilot-servo-command-gates}
\begin{tabularx}{\columnwidth}{Z{0.33\columnwidth}cY}
\toprule
Gate & Value & Meaning \\
\midrule
Manual override block &
\(\SI{0.8}{s}\) &
After an external command, pilot-generated commands are ignored for this long. \\
Command timeout &
\(\SI{1}{s}\) &
A stale nonzero command is reset to zero. \\
Zero-command resend interval &
\(\SI{0.2}{s}\) after \(\SI{1}{s}\) &
Avoids spamming zero commands while still keeping the link alive. \\
Driver timeout &
\(\SI{1}{s}\) &
If command is nonzero but measured current remains zero, a driver-timeout flag is set. \\
\bottomrule
\end{tabularx}
\end{table}
\subsubsection{Current, Voltage, Position, and Rudder Calibration}
The electrical, position-control, and rudder-calibration values are grouped
in \cref{tab:pypilot-servo-electrical-position-rudder} because they are all
low-level actuator calibration parameters rather than steering-loop gains.
\begin{table*}[!tbp]
\centering
\caption{Servo electrical calibration, filtering, position control, and rudder calibration.}
\label{tab:pypilot-servo-electrical-position-rudder}
\begin{tabularx}{\textwidth}{Z{0.23\textwidth}c c Z{0.16\textwidth}Y}
\toprule
Value/filter & Default & Code range & Practical range & Meaning \\
\midrule
\multicolumn{5}{l}{\textit{Servo electrical calibration and filtering}} \\
\addlinespace[0.25ex]
\texttt{servo.current.factor} & \num{1} & \(\num{0.8}\ldots\num{1.2}\) & near calibrated value & Current calibration multiplier. \\
\texttt{servo.current.offset} & \num{0} & \(\num{-1.2}\ldots\num{1.2}\) & near calibrated value & Current calibration offset. \\
\texttt{servo.voltage.factor} & \num{1} & \(\num{0.8}\ldots\num{1.2}\) & near calibrated value & Voltage calibration multiplier. \\
\texttt{servo.voltage.offset} & \num{0} & \(\num{-1.2}\ldots\num{1.2}\) & near calibrated value & Voltage calibration offset. \\
\texttt{servo.compensate\_voltage} & \texttt{false} & boolean & drive-dependent & If enabled, scales command approximately by \(12/V\). \\
\texttt{servo.compensate\_current} & \texttt{false} & boolean & drive-dependent & Uses current compensation when sending controller parameters. \\
Duty EMA alpha & \num{0.001} & fixed & fixed & Very slow estimate of actuator duty cycle. \\
Wattage EMA coefficient & \(0.003\Delta t\) & fixed & fixed & Slow average power estimate, roughly several minutes. \\
\addlinespace[0.6ex]
\multicolumn{5}{l}{\textit{Clutch, brake, and position-control parameters}} \\
\addlinespace[0.25ex]
\texttt{servo.clutch\_pwm} & \(\SI{100}{\percent}\) & \(\SI{10}{\percent}\ldots\SI{100}{\percent}\) & \(\SI{70}{\percent}\ldots\SI{100}{\percent}\) & Clutch PWM command. \\
\texttt{servo.use\_brake} & \texttt{true} & boolean & drive-dependent & Applies brake behavior when idle. \\
\texttt{servo.position.p} & \num{0.1} & \(\num{0.01}\ldots\num{1.0}\) & \(\num{0.05}\ldots\num{0.3}\) & Rudder-position proportional gain. \\
\texttt{servo.position.i} & \num{0} & \(\num{0}\ldots\num{0.1}\) & \(\num{0}\ldots\num{0.02}\) & Rudder-position integral gain. \\
\texttt{servo.position.d} & \num{0.01} & \(\num{0}\ldots\num{0.1}\) & \(\num{0}\ldots\num{0.03}\) & Rudder-position derivative/speed damping. \\
\addlinespace[0.6ex]
\multicolumn{5}{l}{\textit{Rudder calibration values}} \\
\addlinespace[0.25ex]
\texttt{rudder.range} & \(45^\circ\) & \(10^\circ\ldots100^\circ\) & actual hardover angle & Maximum absolute rudder angle. \\
\texttt{rudder.offset} & \num{0} & not range-limited & calibration result & Rudder zero offset. \\
\texttt{rudder.scale} & \num{100} & not range-limited & calibration result & Raw-to-degree scale. \\
\texttt{rudder.nonlinearity} & \num{0} & not range-limited & calibration result & Quadratic correction term. \\
\bottomrule
\end{tabularx}
\end{table*}
Small measured currents are suppressed using the learned noise floor:
\[
I < 1.2 I_{\mathrm{noise}}
\quad\Longrightarrow\quad
I=0.
\]
The current noise estimate is updated only when the command has been idle
for more than \(\SI{3}{s}\), so that motor current is not mistaken for
sensor noise.
The position-command path filters the position error as
\[
e_f[k]
=
0.98\,e_f[k-1]
+
0.02\,\operatorname{clamp}(e[k],-30,30).
\]
At the default control rate, this is a slow filter with an approximate
time constant of \(\SI{2.5}{s}\).  Rudder calibration maps the raw reading
\(r\) to an angle using scale \(s\), offset \(o\), and nonlinearity \(n\):
\[
\delta_r
=
s r
+
o
+
n(r_{\min}-r)(r_{\max}-r).
\]
The rudder-speed estimate is filtered by
\[
\dot{\delta}_{r,f}[k]
=
0.9\,\dot{\delta}_{r,f}[k-1]
+
0.1\,\dot{\delta}_{r}[k].
\]
If the rudder update interval exceeds \(\SI{1}{s}\), the derivative
calculation caps the interval at \(\SI{1}{s}\), preventing a long update
gap from corrupting the speed estimate.
\subsection{Sensor Source Selection and Timeout Gates}
pypilot can receive the same logical sensor from multiple sources.  It
uses a priority table where lower numbers have higher priority.
\begin{table}[!t]
\centering
\caption{Sensor source priorities.}
\label{tab:pypilot-source-priority}
\begin{tabular}{lc}
\toprule
Source & Priority \\
\midrule
\texttt{gpsd} & 1 \\
\texttt{servo} & 1 \\
\texttt{serial} & 2 \\
\texttt{tcp} & 3 \\
\texttt{signalk} & 4 \\
\texttt{water+wind} & 5 \\
\texttt{gps+wind} & 6 \\
\texttt{none} & 7 \\
\bottomrule
\end{tabular}
\end{table}
If a higher-priority source is already active, lower-priority data are
ignored.  If two devices of the same priority provide the same sensor,
pypilot stays with the current device rather than switching randomly.
This prevents source flapping.
A sensor source is dropped if no update has been received for more than
\[
\SI{8}{s}.
\]
The same \(\SI{8}{s}\) stale-value policy is used for servo temperature
sensor values.  A reasonable source-edit range for this timeout is
approximately \(\SI{5}{s}\) to \(\SI{15}{s}\): shorter timeouts react
faster to failures, while longer timeouts tolerate slow or bursty sensors.
\subsection{Mode Availability and Fallback Gates}
The autopilot does not expose all modes at all times.  Available modes are
constructed from the currently available sensors:
\[
\begin{array}{ll}
\text{always available:} & \texttt{compass},\\
\text{if GPS exists:} & \texttt{gps},\\
\text{if APB/navigation exists:} & \texttt{nav},\\
\text{if apparent wind exists:} & \texttt{wind},\\
\text{if true wind exists:} & \texttt{true wind}.
\end{array}
\]
The value
\[
\texttt{ap.gps\_and\_nav\_modes}
\]
defaults to \texttt{true} and controls whether GPS/NAV modes are exposed
when the corresponding sensors exist.
Fallback behavior is summarized in
\cref{tab:pypilot-mode-fallbacks}.
\begin{table}[!t]
\centering
\caption{Mode fallback behavior.}
\label{tab:pypilot-mode-fallbacks}
\begin{tabularx}{\columnwidth}{Z{0.36\columnwidth}Y}
\toprule
Requested mode unavailable & Fallback behavior \\
\midrule
\texttt{nav} &
Falls back toward \texttt{gps}, then \texttt{compass}. \\
\texttt{gps} &
Falls back to \texttt{compass}. \\
\texttt{wind} &
Falls back to \texttt{compass}. \\
\texttt{true wind} &
Falls back to \texttt{wind}. \\
\bottomrule
\end{tabularx}
\end{table}
The wind pilot also performs its own safety fallback: if wind input is
lost, it switches back to the basic pilot.
\subsection{Loop Timing Gates}
The main control period is
\[
T = \frac{1}{\texttt{imu.rate}}.
\]
At the default \(\SI{20}{Hz}\), \(T=\SI{0.05}{s}\).  The main loop
contains diagnostic timing gates, summarized in
\cref{tab:pypilot-loop-timing}.
\begin{table*}[!tbp]
\centering
\caption{Main-loop timing gates.}
\label{tab:pypilot-loop-timing}
\begin{tabularx}{\textwidth}{Z{0.25\textwidth}cZ{0.12\textwidth}Y}
\toprule
Gate & Value & Tunable & Meaning \\
\midrule
Main loop period &
\(1/\texttt{imu.rate}\) &
Yes &
Default control period is \(\SI{0.05}{s}\). \\
IMU read attempts &
14 &
No &
The loop tries multiple times to obtain an IMU reading. \\
Sleep per failed IMU read &
\(T/10\) &
No &
Small retry delay between IMU read attempts. \\
Slow server/client warning &
\(>T/2\) plus standby wait &
No &
Warns if server/client processing is slow. \\
Slow sensor warning &
\(>T/2\) &
No &
Warns if sensor polling is slow. \\
Slow IMU warning &
\(>2T/3\), after startup &
No &
Warns if IMU reading is slow. \\
Slow autopilot warning &
\(>T/2\) &
No &
Warns if control computation is slow. \\
Slow servo warning &
\(>T/2\), when driver exists &
No &
Warns if servo polling is slow. \\
Full iteration warning &
\(>T\), after startup and when enabled &
No &
Warns if the complete loop misses its control period. \\
\bottomrule
\end{tabularx}
\end{table*}
These timing gates are diagnostic rather than control gains.  They are
useful when assessing whether a platform is fast enough to run pypilot
reliably.
\subsection{Disabled or Experimental Pilot Parameters}
The source tree contains additional pilot modules that are disabled by
default.  These include \texttt{simple}, \texttt{rate}, \texttt{gps},
\texttt{deadzone}, \texttt{autotune}, \texttt{fuzzy}, \texttt{vmg}, and
\texttt{learning}.  They define additional gains and parameters, but they
should not be presented as ordinary production tuning parameters unless
the corresponding module is explicitly enabled and tested.

\begin{table*}[!tbp]
\centering
\caption{Disabled or experimental pilot parameters.}
\label{tab:pypilot-disabled-pilots}
\begin{tabularx}{\textwidth}{Z{0.16\textwidth}Z{0.56\textwidth}Z{0.24\textwidth}}
\toprule
Pilot & Representative parameters & Status \\
\midrule
\texttt{simple} &
\texttt{P=0.005}, \texttt{I=0}, \texttt{D=0.15} &
Disabled \\
\texttt{gps} &
\texttt{P=0.003}, \texttt{D=0.1}, \texttt{DD=0.05}, \texttt{FF=0.6} &
Disabled \\
\texttt{rate} &
\texttt{D=0.075}, \texttt{DD=0.075}, \texttt{FF=0.6},
\texttt{maxturnrate=2}, \texttt{turnraterate=0.5} &
Disabled \\
\texttt{deadzone} &
\texttt{deadzone=5}, plus \texttt{P}, \texttt{D}, \texttt{DD}, and
\texttt{FF} gains &
Disabled \\
\texttt{vmg} &
Heading-control gains plus VMG tables &
Disabled/incomplete \\
\texttt{autotune} &
Pilot gains plus cost logic &
Disabled/experimental \\
\texttt{fuzzy} &
\texttt{learningP}, \texttt{learningD}, \texttt{seastate} &
Disabled/experimental \\
\texttt{learning} &
\texttt{P}, \texttt{D}, \texttt{W}, and TensorFlow Lite prediction concept &
Disabled/experimental \\
\bottomrule
\end{tabularx}
\end{table*}

\subsection{Practical Starting Values}
For a typical small sailboat with a tiller or wheel drive, a conservative
starting configuration is
\[
\begin{array}{ll}
\texttt{ap.pilot} &= \texttt{"basic"},\\
\texttt{ap.pilot.basic.P} &= \num{0.003},\\
\texttt{ap.pilot.basic.D} &= \num{0.09},\\
\texttt{ap.pilot.basic.DD} &= \num{0.05}\ldots\num{0.075},\\
\texttt{ap.pilot.basic.PR} &= \num{0.005},\\
\texttt{ap.pilot.basic.FF} &= \num{0.4}\ldots\num{0.8},\\
\texttt{servo.period} &= \SI{0.4}{s},\\
\texttt{servo.gain} &= \num{1},\\
\texttt{servo.speed.min} &= \SI{40}{\percent}\ldots\SI{70}{\percent},\\
\texttt{servo.speed.max} &= \SI{80}{\percent}\ldots\SI{100}{\percent}.
\end{array}
\]
The current limit must be chosen for the actual actuator:
\[
\texttt{servo.max\_current}
=
\text{safe current limit of the drive}.
\]
For wind steering, a conservative starting configuration is
\[
\begin{array}{ll}
\texttt{ap.pilot} &= \texttt{"wind"},\\
\texttt{ap.pilot.wind.P} &= \num{0.003},\\
\texttt{ap.pilot.wind.D} &= \num{0.1},\\
\texttt{ap.pilot.wind.DD} &= \num{0.05},\\
\texttt{ap.pilot.wind.WG} &= \num{0},\\
\texttt{wind.filter\_constant} &= \num{0.05}\ldots\num{0.15},\\
\texttt{ap.wind\_offset\_filter} &= \num{0.05}\ldots\num{0.15}.
\end{array}
\]
For NAV/APB steering, the default cross-track gain is aggressive enough
for many applications:
\[
\texttt{apb.xte.gain}
=
\num{100}\ldots\num{300}
\]
is a reasonable initial range.  Increase it only if the boat returns to
the route line too slowly.  Decrease it if the boat snakes across the
route.
The most safety-critical parameters are
\begin{quote}
\scriptsize\ttfamily
servo.max\_current, servo.max\_controller\_temp, servo.max\_motor\_temp,
rudder.range, servo.speed.max, servo.use\_brake.
\end{quote}
The most important steering-feel parameters are
\begin{quote}
\scriptsize\ttfamily
ap.pilot.basic.P, ap.pilot.basic.D, ap.pilot.basic.PR, ap.pilot.basic.FF,
servo.period, servo.speed.min, servo.speed.max.
\end{quote}

\section{Comparison with Fenix Autopilot}
\label{sec:fenix-comparison}
Fenix Autopilot is an open-source marine tiller-pilot project by Sergio Pascual
(\texttt{spascual90}), first published publicly as an Arduino-based DIY autopilot for
small and medium boats and documented as a self-contained tiller pilot
\cite{fenix-repo,fenix-overview}.  Its architecture and control philosophy differ
substantially from pypilot.  The most important
difference is that pypilot is organized as a distributed autopilot system,
where a Linux/Python process computes pilot commands and communicates with
a separate servo controller, while Fenix is primarily a monolithic Arduino
firmware that performs sensing, state management, control, calibration,
communication, and actuator drive inside the same embedded program.
This distinction affects almost every method in the two systems.  pypilot
is designed around a networked value model, multiple clients, sensor
services, and a separate protected servo layer.  Fenix is designed around a
specific embedded tiller-pilot architecture: an Arduino, an ICM-20948 IMU
or external heading input, a potentiometer-feedback linear actuator, and a
serial/Bluetooth NMEA-style command interface.
\subsection{Architectural Difference}
The pypilot architecture separates the high-level pilot from the actuator
controller.  The Python process computes a normalized servo command, and
the Arduino servo controller converts that command into protected motor
drive.  The servo controller reports current, voltage, temperature, rudder
feedback, and fault flags through a compact binary serial protocol.
Fenix instead combines the autopilot state machine and actuator control in
one firmware image.  Its main control loop parses incoming NMEA and
proprietary Fenix messages, updates heading and sensor state, computes the
pilot command, reads rudder feedback, checks current limits, and directly
drives the linear actuator through PWM and a direction pin.
\begin{table*}[!tbp]
\centering
\caption{Comparison between pypilot and Fenix.}
\label{tab:pypilot-fenix-comparison}
\label{tab:pypilot-fenix-architecture}
\label{tab:pypilot-fenix-control}
\label{tab:pypilot-fenix-actuator}
\label{tab:pypilot-fenix-track}
\label{tab:pypilot-fenix-protocol}
\begingroup
\scriptsize
\setlength{\tabcolsep}{2.4pt}
\renewcommand{\arraystretch}{0.96}
\begin{tabularx}{\textwidth}{Z{0.18\textwidth}YY}
\toprule
Aspect & pypilot & Fenix \\
\midrule
\multicolumn{3}{@{}l}{\textit{System-level architecture}} \\
\midrule
Primary platform &
Linux/Python autopilot process, commonly on Raspberry Pi, with optional
Arduino servo controller. &
Arduino firmware running the complete autopilot and actuator logic. \\
Control distribution &
High-level pilot and low-level servo controller are separate components. &
Heading control, rudder feedback, current protection, calibration, and
actuator drive are integrated in one firmware. \\
Main external interface &
TCP value protocol, web clients, NMEA, Signal K, and serial servo link. &
Serial/Bluetooth NMEA plus proprietary \texttt{\$PEMC} messages. \\
Actuator assumption &
Can drive generic motors through a servo controller; rudder feedback is
optional for the common basic pilot. &
Assumes a linear actuator with potentiometer feedback as a central design
element. \\
Configuration model &
Network-exposed named values, profiles, and runtime-writable properties. &
EEPROM-stored installation parameters, PID gains, and calibration values
set through app/NMEA commissioning messages. \\
\midrule
\multicolumn{3}{@{}l}{\textit{Control-method philosophy}} \\
\midrule
Primary controller output &
Normalized servo/motor command. &
Desired rudder or tiller position. \\
Main active heading law &
Extended PD with heading-rate, heading-acceleration, square-root error,
and feed-forward terms. &
PID heading controller with deadband, derivative filtering, speed-adaptive
gain scaling, integral limiting, and anti-windup. \\
Integral action &
Mostly avoided in active basic and wind pilots. &
Present in the main PID loop, but guarded by anti-windup and actuator
state checks. \\
Derivative action &
Uses measured heading rate and heading rate-of-rate. &
Uses filtered derivative of the heading-error/input signal. \\
Small-error behavior &
Handled by controller gains plus servo pulse/windup logic. &
Handled by explicit heading deadband and rudder feedback deadband. \\
Large-error behavior &
Heading error is clamped before entering the pilot. &
PID output is limited by maximum rudder range. \\
\midrule
\multicolumn{3}{@{}l}{\textit{Actuator control}} \\
\midrule
Actuator command &
Normalized motor effort sent to a servo controller. &
Target rudder position converted to direct PWM/direction drive. \\
Rudder feedback &
Optional in the common basic pilot, but supported by the servo layer. &
Central to normal operation. \\
Low-level motor interface &
Binary serial protocol to an Arduino servo controller. &
Direct Arduino PWM and direction pin. \\
Small movement handling &
Servo-period windup, minimum speed, and pulse shaping. &
Rudder-position deadband, restart threshold, and speed ramping. \\
Direction reversal &
Handled in the servo driver with motor protection logic. &
Explicit stop/delay state before reversing actuator direction. \\
Electrical protection &
Current, voltage, temperature, rudder faults, driver timeout, and
controller flags. &
INA260 overcurrent/blockage detection plus feedback-limit protection. \\
\midrule
\multicolumn{3}{@{}l}{\textit{Track mode}} \\
\midrule
Route input &
APB/GPS data. &
APB data. \\
Internal correction &
Track plus proportional XTE correction. &
Uses external APB course-to-steer. \\
Waypoint management &
Less waypoint-state logic in the pilot. &
Tracks active and next waypoint identifiers and requests confirmation for
new waypoint changes. \\
Main limitation &
XTE law is simple. &
Relies on external navigator for CTS quality. \\
\midrule
\multicolumn{3}{@{}l}{\textit{Commissioning and protocol}} \\
\midrule
High-level control protocol &
TCP value protocol using \texttt{name=json} records. &
NMEA and proprietary \texttt{\$PEMC} sentences over serial/Bluetooth. \\
Low-level actuator protocol &
Separate binary serial protocol to servo controller. &
No separate servo protocol; actuator is driven directly by firmware. \\
Parameter discovery &
Clients can watch the \texttt{values} metadata object. &
Parameters are requested through specific \texttt{\$PEMC} message types. \\
Persistence &
Profiles and persistent values in pypilot configuration. &
EEPROM save commands for installation, PID, and calibration data. \\
Commissioning style &
Networked and service-oriented. &
Embedded marine-instrument style, app/NMEA driven. \\
\bottomrule
\end{tabularx}
\endgroup
\end{table*}
\subsection{Control-Law Philosophy}
The active pypilot basic pilot is an extended proportional--derivative
controller.  Its output is a normalized servo command:
\[
u
=
K_P e
+
K_D \dot{\psi}
+
K_{DD}\ddot{\psi}
+
K_{PR}\operatorname{sgn}(e)\sqrt{\abs{e}}
+
K_{FF}\dot{\psi}_c .
\]
Here \(e\) is heading error, \(\dot{\psi}\) is heading rate,
\(\ddot{\psi}\) is heading acceleration, and \(\dot{\psi}_c\) is command
rate.  The active pypilot basic controller does not rely on an integral
term.  Instead, it uses damping, command feed-forward, a square-root error
term, and motor-side pulse shaping.
Fenix uses a more conventional heading-error PID loop, but with important
extensions.  It computes a target-relative heading error,
\[
e_\psi
=
\operatorname{wrap}_{[-180^\circ,180^\circ)}
\left(
\psi_c-\psi
\right),
\]
then applies an enhanced PID controller.  The PID output is not a raw
motor command.  It is interpreted as a target rudder position:
\[
e_\psi
\longrightarrow
\delta_{r,c}.
\]
A second layer then moves the linear actuator until the measured rudder
position approaches \(\delta_{r,c}\).  This makes Fenix more explicitly a
rudder-position autopilot, while pypilot's basic pilot is more directly a
servo-command autopilot.
\subsection{Speed Adaptation and Deadband}
Fenix contains a speed-adaptive PID scaling method.  A speed factor is
computed from the ratio between a reference speed and the current boat
speed.  The effective gains are then adjusted approximately as
\[
\begin{aligned}
K_{P,\mathrm{eff}} &= K_P f_v,\\
K_{I,\mathrm{eff}} &= K_I f_v,\\
K_{D,\mathrm{eff}} &= \frac{K_D}{f_v}.
\end{aligned}
\]
Thus, when the boat is slower than its nominal reference speed, the
controller increases proportional and integral authority.  When the boat
is faster, proportional and integral gains are reduced and derivative
damping becomes relatively stronger.  This is a practical attempt to keep
rudder response more consistent across different speeds.
Fenix also has an adaptive deadband.  In automatic deadband mode, it
estimates the short-term variation of heading error and chooses a
deadband between a minimum and maximum value:
\[
d_b
=
\operatorname{clamp}
\left(
2 \hat{\sigma}_e,
d_{\min},
d_{\max}
\right).
\]
The default limits are approximately
\[
d_{\min}=0.5^\circ,
\qquad
d_{\max}=5^\circ .
\]
This deadband suppresses unnecessary PID updates near the commanded
heading and reduces actuator hunting.
pypilot does not use this same heading-deadband concept in its active
basic pilot.  Instead, pypilot combines heading-error clamping with the
servo layer's pulse accumulation, minimum movement, and direction-change
logic.  In pypilot, much of the small-command suppression is deliberately
pushed down into the servo driver.
\subsection{Actuator Method}
The actuator strategy is one of the clearest differences.  Fenix assumes a
linear actuator with potentiometer feedback.  The control chain is
\[
\begin{aligned}
e_\psi
&\longrightarrow \delta_{r,c}\\
&\longrightarrow \text{actuator position loop}\\
&\longrightarrow \text{PWM/DIR drive}.
\end{aligned}
\]
The rudder feedback signal is mapped into rudder angle using calibrated
minimum and maximum feedback values, a center offset, and a maximum rudder
range.  Fenix then drives the actuator until the measured rudder position
is close enough to the target position.
Fenix also adds practical actuator logic: it stops at calibrated feedback
limits, uses a rudder-position deadband, delays direction reversals,
ramps motor speed, and reduces speed near the target.  It can detect an
actuator blockage using an INA260 current monitor and stop the drive when
current exceeds the configured limit.
pypilot's common active pilot does not require that the actuator position
be known.  Its servo layer converts normalized motor commands into motor
pulses and protects the drive using current, voltage, temperature, rudder
feedback when available, driver timeout, and fault flags from the Arduino
servo controller.  This makes pypilot more tolerant of different actuator
types, while Fenix is more focused on a known feedback linear actuator.
\subsection{Heading and Sensor Handling}
Fenix uses an ICM-20948 IMU with a Mahony AHRS implementation for internal
heading.  It can also accept external heading through NMEA heading
messages.  When external heading is active, Fenix extrapolates heading
between received messages using the most recent heading slope.  This helps
keep heading continuous when external heading messages arrive at a lower
rate than the control loop.
pypilot is more sensor-service oriented.  Its internal IMU path, external
sensor paths, and networked values are integrated into a larger runtime
environment.  pypilot is therefore better suited to systems where heading,
wind, GPS, rudder, and other values may come from multiple network or
serial sources.
Fenix includes validity timeouts for external heading, speed, wind, and
APB/navigation data.  If external heading is lost while operating from
that source, it can return to standby and raise an alarm.  pypilot has
similar safety behavior at the system level, but the implementation is
distributed across the sensor/value framework and autopilot mode logic.
\subsection{Navigation and Track Mode}
Both systems support route following, but neither implements a full modern
path-following controller such as L1, pure pursuit, or Stanley control in
the inspected active code.
pypilot's NAV mode computes a heading command using a simple cross-track
correction:
\[
\psi_c
=
\psi_{\mathrm{track}}
+
K_{\mathrm{xte}} e_{\mathrm{xte}} .
\]
Thus pypilot performs its own proportional cross-track-error correction.
Fenix instead receives APB navigation data and uses the externally
supplied course-to-steer value:
\[
\psi_c = \psi_{\mathrm{CTS}} .
\]
The APB sentence also contains cross-track error, bearings, alarms, and
waypoint identifiers, but the active Fenix track-mode steering path uses
the APB course-to-steer field as the target bearing.  In practice, this
delegates more of the navigation calculation to the plotter or navigation
source.
\subsection{Wind Mode}
pypilot has an active wind pilot with its own gains.  It uses heading
error, heading-rate damping, heading-acceleration damping, and optionally
a wind-gust term.  The wind pilot therefore has a separate steering law
from the basic compass pilot.
Fenix implements wind steering by continuously converting wind-angle hold
into a target bearing.  It supports both apparent-wind and true-wind
inputs.  If \(\theta_w\) is the current wind-relative direction and
\(\theta_{w,c}\) is the target wind-relative direction, Fenix updates the
course target approximately as
\[
\psi_c
=
\psi
+
\theta_w
-
\theta_{w,c}.
\]
After this target bearing is computed, the same main heading PID loop
drives the rudder.  This makes the wind mode simpler and more unified, but
less specialized than pypilot's wind-specific controller.
\subsection{Commissioning and Configuration}
Fenix has a strong embedded commissioning model.  Installation parameters,
PID gains, feedback calibration, IMU calibration, and EEPROM saving are
all handled through the app and through proprietary NMEA-style
\texttt{\$PEMC} messages.  These messages can set installation parameters,
set PID gains, enter or exit feedback calibration, load external
ICM-20948 calibration matrices, and save results to EEPROM.
pypilot's commissioning is more dynamic.  Most parameters are exposed as
named values over the TCP value protocol and can be watched or changed by
clients.  Persistent profiles and runtime values are part of the pypilot
server model.  The servo controller also has its own binary protocol for
low-level motor parameters.
\subsection{Safety and Alarms}
Fenix contains explicit marine-style alarms.  It raises alarms for
out-of-course operation, lost external heading, wind changes, and actuator
blockage.  Off-course and wind-change alarms are delayed so that large
deliberate turns do not immediately trigger an alarm.  Fenix also refuses
some parameter changes unless the autopilot is in standby, which is a
useful commissioning safety rule.
pypilot's safety model is more distributed.  The autopilot logic handles
mode fallbacks and sensor availability, while the servo layer handles
motor and electrical protection.  The Arduino servo controller reports
faults such as overcurrent, overtemperature, bad voltage, rudder limits,
and driver status.  pypilot is therefore stronger at low-level electrical
and motor-controller supervision, while Fenix is stronger at integrating
linear-actuator position, calibration, and user-facing commissioning into
one firmware.
\subsection{Practical Interpretation}
Fenix and pypilot represent two different design philosophies.
Fenix is best understood as a self-contained embedded tiller pilot:
{\small
\[
\begin{gathered}
\text{Arduino}+\text{ICM-20948}+{}\\
\text{feedback linear actuator}+\text{NMEA/Bluetooth app}.
\end{gathered}
\]
}
Its methods are well aligned with that hardware model.  It uses a
position-feedback actuator, a PID heading controller, speed-adaptive gain
scaling, explicit deadband, EEPROM commissioning, and direct PWM actuator
drive.  If the actuator feedback is reliable and the boat matches the
intended installation, Fenix's rudder-position approach is mechanically
clear and practical.
pypilot is best understood as an open marine autopilot system:
{\small
\[
\begin{gathered}
\text{Linux process}+\text{network value server}+{}\\
\text{sensor services}+\text{servo controller}.
\end{gathered}
\]
}
Its methods are more extensible and more network-oriented.  It can support
multiple clients, remote displays, web interfaces, Signal K integration,
NMEA bridges, and different pilot algorithms.  Its basic active pilot is
less conventional than Fenix's PID, but its servo controller is more
general-purpose and more richly supervised.
The central method difference can be summarized as

{\small
\[
\begin{array}{ll}
\text{pypilot:} &
\begin{gathered}[t]
\text{heading error}\to\text{motor command}\to{}\\
\text{protected servo layer},
\end{gathered}
\\[6pt]
\text{Fenix:} &
\begin{gathered}[t]
\text{heading error}\to\text{rudder angle}\to{}\\
\text{linear-actuator loop}.
\end{gathered}
\end{array}
\]
}

For a small Arduino-only autopilot with a known potentiometer-feedback
linear actuator, Fenix is a coherent and practical design.  For a larger
open navigation system with network clients, mixed sensors, remote
interfaces, and replaceable pilot algorithms, pypilot has the stronger
system architecture.

\section{Conclusion}

The uploaded \texttt{pypilot} source is best characterized as Sean D'Epagnier's
practical, open, marine heading autopilot rather than a full modern autonomous-navigation
stack \cite{pypilot-author}.  Its default
method is an extended PD heading controller with gyro damping, yaw-acceleration damping,
square-root heading-error shaping, and heading-command feed-forward.  Its strongest
engineering feature is the servo/motor layer, which treats real rudder drives as nonlinear,
limited, fault-prone devices rather than ideal actuators.

Compared with ArduPilot, the community project rooted in Jordi Mu\~noz and Chris
Anderson's early DIY Drones work, or PX4, \texttt{pypilot} has a much simpler guidance
layer and a less formal state-estimation/control hierarchy.  Compared with Fenix
Autopilot, Sergio Pascual's self-contained Arduino tiller-pilot design, \texttt{pypilot}
is more networked and modular, while Fenix gives useful credit to the opposite embedded
philosophy: integrating heading, actuator, calibration, and alarms in one compact unit.
Compared with commercial sailing pilots,
it is less polished and contains incomplete experimental code, but it is transparent,
modifiable, and deeply oriented toward real DIY marine drives.  The most valuable next step
is not to replace it with machine learning, but to add a stronger deterministic guidance
layer above the existing heading and servo architecture.

\FloatBarrier

\section*{Acknowledgments}
This work was prepared with the assistance of artificial intelligence tools, including ChatGPT, Codex.

\begin{thebibliography}{9}

\bibitem{pypilot-author}
pypilot.org,
\emph{Author}.  Credits Sean D'Epagnier as author and open-source marine software contributor.
\url{https://pypilot.org/author/}

\bibitem{fenix-repo}
Sergio Pascual (\texttt{spascual90}),
\emph{Fenix Autopilot}.  Public GitHub repository for the open-source Arduino-based DIY tiller pilot.
\url{https://github.com/spascual90/Fenix}

\bibitem{fenix-overview}
Fenix Autopilot Documentation,
\emph{Overview}.  Describes Fenix as a self-contained tiller pilot for tiller-steered sailboats.
\url{https://spascual90.gitbook.io/fenix-autopilot/introduction/overview}

\bibitem{ardupilot-history}
ArduPilot Documentation,
\emph{History of ArduPilot}.  Summarizes the project's DIY Drones origins and later community development.
\url{https://ardupilot.org/copter/docs/common-history-of-ardupilot.html}

\bibitem{ardupilot-l1}
ArduPilot Dev Team,
\emph{Rover: Tuning Navigation}.  Describes Rover navigation tuning and the L1 controller.
\url{https://ardupilot.org/rover/docs/rover-tuning-navigation-420.html}

\bibitem{ardupilot-rate}
ArduPilot Dev Team,
\emph{Rover: Tuning Turn Rate}.  Describes the steering-rate PID controller and FF/P/I/D tuning.
\url{https://ardupilot.org/rover/docs/rover-tuning-steering-rate.html}

\bibitem{px4-fw}
PX4 Autopilot,
\emph{Fixed-wing Altitude/Position Controller Tuning}.  Describes TECS and NPFG in PX4 fixed-wing control.
\url{https://docs.px4.io/main/en/config_fw/position_tuning_guide_fixedwing}

\bibitem{pypilot-user-manual}
pypilot.org,
\emph{pypilot User Manual}.  User-facing reference for pypilot setup, operation, calibration,
clients, and autopilot modes.
\url{https://pypilot.org/doc/pypilot_user_manual/pdf/pypilot_user_manual.pdf}

\bibitem{pypilot-workbook-wiki}
pypilot project,
\emph{Workbook Wiki}.  Community-maintained operational and commissioning notes for pypilot
hardware, software, calibration, and integration.
\url{https://github.com/pypilot/workbook/tree/main/wiki}

\bibitem{bg-pilot}
B\&G Sailing,
\emph{Autopilot Displays/Controllers}.  Describes Wind and Navigation auto-steering modes and pilot controls.
\url{https://www.bandg.com/en-sg/bg/type/autopilots/autopilot-displays-controllers/}

\bibitem{bg-setup}
B\&G Sailing,
\emph{Autopilots: Setting up your Autopilot}.  Discusses true-wind-angle use in shorthanded gybing practice.
\url{https://www.bandg.com/en-sg/blog/autopilots-setting-up-your-autopilot/}

\end{thebibliography}

\end{document}
